% =====================================================================
%  Anatomía de un LLM — Fase 0 — Módulo 0.7
%  Vectores y funciones de varias variables
%  Compilar: latexmk -r ../latexmkrc P07-varias-variables.tex
% =====================================================================
\documentclass[11pt,a4paper,twoside,openany]{book}
\usepackage{../estilo-apuntes}

\newcommand{\ApuntesCurso}{Anatomía de un LLM}
\newcommand{\ApuntesModulo}{Módulo 0.7}

% Anclas de hiperenlace únicas por tema para los enunciados
\makeatletter
\@for\nombre:=definicion,teorema,proposicion,lema,corolario,ejemplo,observacion\do{%
  \expandafter\xdef\csname theH\nombre\endcsname{\noexpand\thechapter.\noexpand\arabic{definicion}}}
\makeatother

\newcommand{\vu}{\mathbf{u}}
\newcommand{\vv}{\mathbf{v}}
\newcommand{\vw}{\mathbf{w}}
\newcommand{\vn}{\mathbf{n}}
\newcommand{\vx}{\mathbf{x}}
\newcommand{\va}{\mathbf{a}}
\newcommand{\vb}{\mathbf{b}}
\newcommand{\vr}{\mathbf{r}}
\newcommand{\vi}{\mathbf{i}}
\newcommand{\vj}{\mathbf{j}}
\newcommand{\vk}{\mathbf{k}}
\newcommand{\vcero}{\mathbf{0}}
\newcommand{\Z}{\mathbb{Z}}
\newcommand{\proj}{\operatorname{proj}}
\newcommand{\Dom}{\operatorname{Dom}}
\newcommand{\dA}{\,dA}
\newcommand{\pd}[2]{\frac{\partial #1}{\partial #2}}

\hypersetup{pdftitle={Módulo 0.7. Vectores y funciones de varias variables (resumen)},pdfauthor={Anatomía de un LLM}}

\makeatletter
\newlabel{ej:t1-punto-medio}{{1.2}{1}{}{ejemplo.1.2}{}}
\newlabel{ej:t1-punto-medio@cref}{{[ejemplo][2][]1.2}{[1][1][]1}}
\newlabel{ej:t1-angulo}{{1.9}{3}{}{ejemplo.1.9}{}}
\newlabel{ej:t1-angulo@cref}{{[ejemplo][9][]1.9}{[1][3][]3}}
\newlabel{fig:t1-proyeccion}{{1.1}{3}{Proyección ortogonal de $\vu $ sobre $\vv $. Su longitud con signo es $\norm {\vu }\cos \theta =\operatorname {comp}_{\vv }\vu $}{figure.1.1}{}}
\newlabel{fig:t1-proyeccion@cref}{{[figure][1][1]1.1}{[1][3][]3}}
\newlabel{ej:t1-proyeccion}{{1.12}{3}{}{ejemplo.1.12}{}}
\newlabel{ej:t1-proyeccion@cref}{{[ejemplo][12][]1.12}{[1][3][]3}}
\newlabel{ej:t1-volumen}{{1.17}{4}{}{ejemplo.1.17}{}}
\newlabel{ej:t1-volumen@cref}{{[ejemplo][17][]1.17}{[1][4][]4}}
\newlabel{eje:1-vf}{{1.1}{5}{}{ejercicio.1.1}{}}
\newlabel{eje:1-vf@cref}{{[ejercicio][1][1]1.1}{[1][5][]5}}
\newlabel{eje:1-cosenos}{{1.2}{5}{}{ejercicio.1.2}{}}
\newlabel{eje:1-cosenos@cref}{{[ejercicio][2][1]1.2}{[1][5][]5}}
\newlabel{eje:1-operaciones}{{1.3}{5}{}{ejercicio.1.3}{}}
\newlabel{eje:1-operaciones@cref}{{[ejercicio][3][1]1.3}{[1][5][]5}}
\newlabel{eje:1-unitarios}{{1.4}{5}{}{ejercicio.1.4}{}}
\newlabel{eje:1-unitarios@cref}{{[ejercicio][4][1]1.4}{[1][5][]5}}
\newlabel{eje:1-angulos}{{1.5}{5}{}{ejercicio.1.5}{}}
\newlabel{eje:1-angulos@cref}{{[ejercicio][5][1]1.5}{[1][5][]5}}
\newlabel{eje:1-ortogonal-c}{{1.6}{5}{}{ejercicio.1.6}{}}
\newlabel{eje:1-ortogonal-c@cref}{{[ejercicio][6][1]1.6}{[1][5][]5}}
\newlabel{eje:1-proyeccion}{{1.7}{6}{}{ejercicio.1.7}{}}
\newlabel{eje:1-proyeccion@cref}{{[ejercicio][7][1]1.7}{[1][5][]6}}
\newlabel{eje:1-cubo}{{1.8}{6}{}{ejercicio.1.8}{}}
\newlabel{eje:1-cubo@cref}{{[ejercicio][8][1]1.8}{[1][6][]6}}
\newlabel{eje:1-vectorial}{{1.9}{6}{}{ejercicio.1.9}{}}
\newlabel{eje:1-vectorial@cref}{{[ejercicio][9][1]1.9}{[1][6][]6}}
\newlabel{eje:1-tetraedro}{{1.10}{6}{}{ejercicio.1.10}{}}
\newlabel{eje:1-tetraedro@cref}{{[ejercicio][10][1]1.10}{[1][6][]6}}
\newlabel{eje:1-polarizacion}{{1.11}{6}{}{ejercicio.1.11}{}}
\newlabel{eje:1-polarizacion@cref}{{[ejercicio][11][1]1.11}{[1][6][]6}}
\newlabel{eje:1-cs-cuadratica}{{1.12}{6}{}{ejercicio.1.12}{}}
\newlabel{eje:1-cs-cuadratica@cref}{{[ejercicio][12][1]1.12}{[1][6][]6}}
\newlabel{eje:1-rombo}{{1.13}{6}{}{ejercicio.1.13}{}}
\newlabel{eje:1-rombo@cref}{{[ejercicio][13][1]1.13}{[1][6][]6}}
\newlabel{eje:1-semicirculo}{{1.14}{6}{}{ejercicio.1.14}{}}
\newlabel{eje:1-semicirculo@cref}{{[ejercicio][14][1]1.14}{[1][6][]6}}
\newlabel{eje:1-lagrange-id}{{1.15}{6}{}{ejercicio.1.15}{}}
\newlabel{eje:1-lagrange-id@cref}{{[ejercicio][15][1]1.15}{[1][6][]6}}
\newlabel{ej:t2-plano}{{2.3}{7}{}{ejemplo.2.3}{}}
\newlabel{ej:t2-plano@cref}{{[ejemplo][3][]2.3}{[1][7][]7}}
\newlabel{ej:t2-recta-interseccion}{{2.5}{8}{}{ejemplo.2.5}{}}
\newlabel{ej:t2-recta-interseccion@cref}{{[ejemplo][5][]2.5}{[1][8][]8}}
\newlabel{ej:t2-gauss}{{2.9}{9}{}{ejemplo.2.9}{}}
\newlabel{ej:t2-gauss@cref}{{[ejemplo][9][]2.9}{[1][9][]9}}
\newlabel{eje:2-posicion}{{2.1}{11}{}{ejercicio.2.1}{}}
\newlabel{eje:2-posicion@cref}{{[ejercicio][1][2]2.1}{[1][11][]11}}
\newlabel{eje:2-planos}{{2.2}{11}{}{ejercicio.2.2}{}}
\newlabel{eje:2-planos@cref}{{[ejercicio][2][2]2.2}{[1][11][]11}}
\newlabel{eje:2-diedro}{{2.3}{11}{}{ejercicio.2.3}{}}
\newlabel{eje:2-diedro@cref}{{[ejercicio][3][2]2.3}{[1][11][]11}}
\newlabel{eje:2-interseccion}{{2.4}{11}{}{ejercicio.2.4}{}}
\newlabel{eje:2-interseccion@cref}{{[ejercicio][4][2]2.4}{[1][11][]11}}
\newlabel{eje:2-gauss}{{2.5}{11}{}{ejercicio.2.5}{}}
\newlabel{eje:2-gauss@cref}{{[ejercicio][5][2]2.5}{[1][11][]11}}
\newlabel{eje:2-gauss-casos}{{2.6}{11}{}{ejercicio.2.6}{}}
\newlabel{eje:2-gauss-casos@cref}{{[ejercicio][6][2]2.6}{[1][11][]11}}
\newlabel{eje:2-parametro}{{2.7}{11}{}{ejercicio.2.7}{}}
\newlabel{eje:2-parametro@cref}{{[ejercicio][7][2]2.7}{[1][11][]11}}
\newlabel{eje:2-areas}{{2.8}{11}{}{ejercicio.2.8}{}}
\newlabel{eje:2-areas@cref}{{[ejercicio][8][2]2.8}{[1][11][]11}}
\newlabel{eje:2-inversa}{{2.9}{11}{}{ejercicio.2.9}{}}
\newlabel{eje:2-inversa@cref}{{[ejercicio][9][2]2.9}{[1][11][]11}}
\newlabel{eje:2-produccion}{{2.10}{12}{}{ejercicio.2.10}{}}
\newlabel{eje:2-produccion@cref}{{[ejercicio][10][2]2.10}{[1][11][]12}}
\newlabel{eje:2-cramer}{{2.11}{12}{}{ejercicio.2.11}{}}
\newlabel{eje:2-cramer@cref}{{[ejercicio][11][2]2.11}{[1][12][]12}}
\newlabel{eje:2-distancia}{{2.12}{12}{}{ejercicio.2.12}{}}
\newlabel{eje:2-distancia@cref}{{[ejercicio][12][2]2.12}{[1][12][]12}}
\newlabel{eje:2-vandermonde}{{2.13}{12}{}{ejercicio.2.13}{}}
\newlabel{eje:2-vandermonde@cref}{{[ejercicio][13][2]2.13}{[1][12][]12}}
\newlabel{eje:2-compatibilidad}{{2.14}{12}{}{ejercicio.2.14}{}}
\newlabel{eje:2-compatibilidad@cref}{{[ejercicio][14][2]2.14}{[1][12][]12}}
\newlabel{eje:2-producto-det}{{2.15}{12}{}{ejercicio.2.15}{}}
\newlabel{eje:2-producto-det@cref}{{[ejercicio][15][2]2.15}{[1][12][]12}}
\newlabel{ej:t3-graficas}{{3.2}{13}{}{ejemplo.3.2}{}}
\newlabel{ej:t3-graficas@cref}{{[ejemplo][2][]3.2}{[1][13][]13}}
\newlabel{ej:t3-niveles}{{3.4}{13}{}{ejemplo.3.4}{}}
\newlabel{ej:t3-niveles@cref}{{[ejemplo][4][]3.4}{[1][13][]13}}
\newlabel{fig:t3-niveles}{{3.1}{14}{Curvas de nivel del paraboloide $x^2+y^2$ (circunferencias, cada vez más juntas porque la función crece cada vez más deprisa) y de la silla $x^2-y^2$ (hipérbolas: trazo continuo para $c>0$, discontinuo para $c<0$; el nivel~$0$ son las dos diagonales)}{figure.3.1}{}}
\newlabel{fig:t3-niveles@cref}{{[figure][1][3]3.1}{[1][13][]14}}
\newlabel{ej:t3-polares}{{3.7}{15}{}{ejemplo.3.7}{}}
\newlabel{ej:t3-polares@cref}{{[ejemplo][7][]3.7}{[1][14][]15}}
\newlabel{ej:t3-no-existe}{{3.10}{15}{}{ejemplo.3.10}{}}
\newlabel{ej:t3-no-existe@cref}{{[ejemplo][10][]3.10}{[1][15][]15}}
\newlabel{ej:t3-trozos}{{3.13}{15}{}{ejemplo.3.13}{}}
\newlabel{ej:t3-trozos@cref}{{[ejemplo][13][]3.13}{[1][15][]15}}
\newlabel{eje:3-conceptos}{{3.1}{16}{}{ejercicio.3.1}{}}
\newlabel{eje:3-conceptos@cref}{{[ejercicio][1][3]3.1}{[1][16][]16}}
\newlabel{eje:3-dominios}{{3.2}{16}{}{ejercicio.3.2}{}}
\newlabel{eje:3-dominios@cref}{{[ejercicio][2][3]3.2}{[1][16][]16}}
\newlabel{eje:3-niveles}{{3.3}{16}{}{ejercicio.3.3}{}}
\newlabel{eje:3-niveles@cref}{{[ejercicio][3][3]3.3}{[1][16][]16}}
\newlabel{eje:3-niveles-2}{{3.4}{16}{}{ejercicio.3.4}{}}
\newlabel{eje:3-niveles-2@cref}{{[ejercicio][4][3]3.4}{[1][16][]16}}
\newlabel{eje:3-onda}{{3.5}{16}{}{ejercicio.3.5}{}}
\newlabel{eje:3-onda@cref}{{[ejercicio][5][3]3.5}{[1][16][]16}}
\newlabel{eje:3-interseccion-superficies}{{3.6}{17}{}{ejercicio.3.6}{}}
\newlabel{eje:3-interseccion-superficies@cref}{{[ejercicio][6][3]3.6}{[1][17][]17}}
\newlabel{eje:3-limites-directos}{{3.7}{17}{}{ejercicio.3.7}{}}
\newlabel{eje:3-limites-directos@cref}{{[ejercicio][7][3]3.7}{[1][17][]17}}
\newlabel{eje:3-no-existe}{{3.8}{17}{}{ejercicio.3.8}{}}
\newlabel{eje:3-no-existe@cref}{{[ejercicio][8][3]3.8}{[1][17][]17}}
\newlabel{eje:3-polares}{{3.9}{17}{}{ejercicio.3.9}{}}
\newlabel{eje:3-polares@cref}{{[ejercicio][9][3]3.9}{[1][17][]17}}
\newlabel{eje:3-continuidad}{{3.10}{17}{}{ejercicio.3.10}{}}
\newlabel{eje:3-continuidad@cref}{{[ejercicio][10][3]3.10}{[1][17][]17}}
\newlabel{eje:3-lse}{{3.11}{17}{}{ejercicio.3.11}{}}
\newlabel{eje:3-lse@cref}{{[ejercicio][11][3]3.11}{[1][17][]17}}
\newlabel{eje:3-epsilon-delta}{{3.12}{17}{}{ejercicio.3.12}{}}
\newlabel{eje:3-epsilon-delta@cref}{{[ejercicio][12][3]3.12}{[1][17][]17}}
\newlabel{eje:3-composicion}{{3.13}{17}{}{ejercicio.3.13}{}}
\newlabel{eje:3-composicion@cref}{{[ejercicio][13][3]3.13}{[1][17][]17}}
\newlabel{eje:3-cuadratica}{{3.14}{17}{}{ejercicio.3.14}{}}
\newlabel{eje:3-cuadratica@cref}{{[ejercicio][14][3]3.14}{[1][17][]17}}
\newlabel{eje:3-trayectorias}{{3.15}{18}{}{ejercicio.3.15}{}}
\newlabel{eje:3-trayectorias@cref}{{[ejercicio][15][3]3.15}{[1][17][]18}}
\newlabel{ej:t4-parciales}{{4.2}{19}{}{ejemplo.4.2}{}}
\newlabel{ej:t4-parciales@cref}{{[ejemplo][2][]4.2}{[1][19][]19}}
\newlabel{ej:t4-aproximacion}{{4.7}{20}{}{ejemplo.4.7}{}}
\newlabel{ej:t4-aproximacion@cref}{{[ejemplo][7][]4.7}{[1][20][]20}}
\newlabel{fig:t4-arbol}{{4.1}{21}{Diagrama de árbol de $w=f(x,y)$ con $x=x(u,v)$, $y=y(u,v)$. Los dos caminos de $w$ a $u$ dan $w_u=w_xx_u+w_yy_u$}{figure.4.1}{}}
\newlabel{fig:t4-arbol@cref}{{[figure][1][4]4.1}{[1][21][]21}}
\newlabel{ej:t4-polares}{{4.9}{21}{}{ejemplo.4.9}{}}
\newlabel{ej:t4-polares@cref}{{[ejemplo][9][]4.9}{[1][21][]21}}
\newlabel{ej:t4-direccional}{{4.14}{22}{}{ejemplo.4.14}{}}
\newlabel{ej:t4-direccional@cref}{{[ejemplo][14][]4.14}{[1][22][]22}}
\newlabel{eje:4-conceptos}{{4.1}{23}{}{ejercicio.4.1}{}}
\newlabel{eje:4-conceptos@cref}{{[ejercicio][1][4]4.1}{[1][22][]23}}
\newlabel{eje:4-parciales}{{4.2}{23}{}{ejercicio.4.2}{}}
\newlabel{eje:4-parciales@cref}{{[ejercicio][2][4]4.2}{[1][23][]23}}
\newlabel{eje:4-schwarz}{{4.3}{23}{}{ejercicio.4.3}{}}
\newlabel{eje:4-schwarz@cref}{{[ejercicio][3][4]4.3}{[1][23][]23}}
\newlabel{eje:4-tangente}{{4.4}{23}{}{ejercicio.4.4}{}}
\newlabel{eje:4-tangente@cref}{{[ejercicio][4][4]4.4}{[1][23][]23}}
\newlabel{eje:4-aproximacion}{{4.5}{23}{}{ejercicio.4.5}{}}
\newlabel{eje:4-aproximacion@cref}{{[ejercicio][5][4]4.5}{[1][23][]23}}
\newlabel{eje:4-gradiente}{{4.6}{23}{}{ejercicio.4.6}{}}
\newlabel{eje:4-gradiente@cref}{{[ejercicio][6][4]4.6}{[1][23][]23}}
\newlabel{eje:4-max-direccional}{{4.7}{23}{}{ejercicio.4.7}{}}
\newlabel{eje:4-max-direccional@cref}{{[ejercicio][7][4]4.7}{[1][23][]23}}
\newlabel{eje:4-tangente-nivel}{{4.8}{23}{}{ejercicio.4.8}{}}
\newlabel{eje:4-tangente-nivel@cref}{{[ejercicio][8][4]4.8}{[1][23][]23}}
\newlabel{eje:4-cadena-t}{{4.9}{23}{}{ejercicio.4.9}{}}
\newlabel{eje:4-cadena-t@cref}{{[ejercicio][9][4]4.9}{[1][23][]23}}
\newlabel{eje:4-cadena-uv}{{4.10}{23}{}{ejercicio.4.10}{}}
\newlabel{eje:4-cadena-uv@cref}{{[ejercicio][10][4]4.10}{[1][23][]23}}
\newlabel{eje:4-softmax}{{4.11}{24}{}{ejercicio.4.11}{}}
\newlabel{eje:4-softmax@cref}{{[ejercicio][11][4]4.11}{[1][23][]24}}
\newlabel{eje:4-polares}{{4.12}{24}{}{ejercicio.4.12}{}}
\newlabel{eje:4-polares@cref}{{[ejercicio][12][4]4.12}{[1][24][]24}}
\newlabel{eje:4-euler}{{4.13}{24}{}{ejercicio.4.13}{}}
\newlabel{eje:4-euler@cref}{{[ejercicio][13][4]4.13}{[1][24][]24}}
\newlabel{eje:4-material}{{4.14}{24}{}{ejercicio.4.14}{}}
\newlabel{eje:4-material@cref}{{[ejercicio][14][4]4.14}{[1][24][]24}}
\newlabel{eje:4-no-diferenciable}{{4.15}{24}{}{ejercicio.4.15}{}}
\newlabel{eje:4-no-diferenciable@cref}{{[ejercicio][15][4]4.15}{[1][24][]24}}
\newlabel{ej:t5-hessiano}{{5.6}{26}{}{ejemplo.5.6}{}}
\newlabel{ej:t5-hessiano@cref}{{[ejemplo][6][]5.6}{[1][26][]26}}
\newlabel{ej:t5-absolutos}{{5.8}{27}{}{ejemplo.5.8}{}}
\newlabel{ej:t5-absolutos@cref}{{[ejemplo][8][]5.8}{[1][27][]27}}
\newlabel{fig:t5-lagrange}{{5.1}{27}{Máximo de $f(x,y)=x+y$ sobre la circunferencia $g(x,y)=x^2+y^2=2$. Las rectas de nivel de $f$ (discontinuas) cortan a la circunferencia salvo la de nivel $2$, que es tangente en $(1,1)$; allí $\nabla f=(1,1)$ y $\nabla g=(2,2)$ son paralelos}{figure.5.1}{}}
\newlabel{fig:t5-lagrange@cref}{{[figure][1][5]5.1}{[1][27][]27}}
\newlabel{ej:t5-lagrange-circulo}{{5.10}{28}{}{ejemplo.5.10}{}}
\newlabel{ej:t5-lagrange-circulo@cref}{{[ejemplo][10][]5.10}{[1][28][]28}}
\newlabel{ej:t5-lagrange-plano}{{5.11}{28}{}{ejemplo.5.11}{}}
\newlabel{ej:t5-lagrange-plano@cref}{{[ejemplo][11][]5.11}{[1][28][]28}}
\newlabel{eje:5-conceptos}{{5.1}{29}{}{ejercicio.5.1}{}}
\newlabel{eje:5-conceptos@cref}{{[ejercicio][1][5]5.1}{[1][28][]29}}
\newlabel{eje:5-clasificar-1}{{5.2}{29}{}{ejercicio.5.2}{}}
\newlabel{eje:5-clasificar-1@cref}{{[ejercicio][2][5]5.2}{[1][29][]29}}
\newlabel{eje:5-clasificar-2}{{5.3}{29}{}{ejercicio.5.3}{}}
\newlabel{eje:5-clasificar-2@cref}{{[ejercicio][3][5]5.3}{[1][29][]29}}
\newlabel{eje:5-clasificar-3}{{5.4}{29}{}{ejercicio.5.4}{}}
\newlabel{eje:5-clasificar-3@cref}{{[ejercicio][4][5]5.4}{[1][29][]29}}
\newlabel{eje:5-openstax}{{5.5}{29}{}{ejercicio.5.5}{}}
\newlabel{eje:5-openstax@cref}{{[ejercicio][5][5]5.5}{[1][29][]29}}
\newlabel{eje:5-absolutos}{{5.6}{29}{}{ejercicio.5.6}{}}
\newlabel{eje:5-absolutos@cref}{{[ejercicio][6][5]5.6}{[1][29][]29}}
\newlabel{eje:5-disco}{{5.7}{29}{}{ejercicio.5.7}{}}
\newlabel{eje:5-disco@cref}{{[ejercicio][7][5]5.7}{[1][29][]29}}
\newlabel{eje:5-caja}{{5.8}{29}{}{ejercicio.5.8}{}}
\newlabel{eje:5-caja@cref}{{[ejercicio][8][5]5.8}{[1][29][]29}}
\newlabel{eje:5-hiperbola}{{5.9}{29}{}{ejercicio.5.9}{}}
\newlabel{eje:5-hiperbola@cref}{{[ejercicio][9][5]5.9}{[1][29][]29}}
\newlabel{eje:5-resistencias}{{5.10}{29}{}{ejercicio.5.10}{}}
\newlabel{eje:5-resistencias@cref}{{[ejercicio][10][5]5.10}{[1][29][]29}}
\newlabel{eje:5-minimos-cuadrados}{{5.11}{29}{}{ejercicio.5.11}{}}
\newlabel{eje:5-minimos-cuadrados@cref}{{[ejercicio][11][5]5.11}{[1][29][]29}}
\newlabel{eje:5-degenerado}{{5.12}{30}{}{ejercicio.5.12}{}}
\newlabel{eje:5-degenerado@cref}{{[ejercicio][12][5]5.12}{[1][30][]30}}
\newlabel{eje:5-mle}{{5.13}{30}{}{ejercicio.5.13}{}}
\newlabel{eje:5-mle@cref}{{[ejercicio][13][5]5.13}{[1][30][]30}}
\newlabel{eje:5-media-geometrica}{{5.14}{30}{}{ejercicio.5.14}{}}
\newlabel{eje:5-media-geometrica@cref}{{[ejercicio][14][5]5.14}{[1][30][]30}}
\newlabel{eje:5-cs-lagrange}{{5.15}{30}{}{ejercicio.5.15}{}}
\newlabel{eje:5-cs-lagrange@cref}{{[ejercicio][15][5]5.15}{[1][30][]30}}
\newlabel{ej:t6-rectangulo}{{6.5}{32}{}{ejemplo.6.5}{}}
\newlabel{ej:t6-rectangulo@cref}{{[ejemplo][5][]6.5}{[1][32][]32}}
\newlabel{ej:t6-region}{{6.7}{32}{}{ejemplo.6.7}{}}
\newlabel{ej:t6-region@cref}{{[ejemplo][7][]6.7}{[1][32][]32}}
\newlabel{ej:t6-orden}{{6.8}{32}{}{ejemplo.6.8}{}}
\newlabel{ej:t6-orden@cref}{{[ejemplo][8][]6.8}{[1][32][]32}}
\newlabel{ej:t6-polares}{{6.10}{33}{}{ejemplo.6.10}{}}
\newlabel{ej:t6-polares@cref}{{[ejemplo][10][]6.10}{[1][33][]33}}
\newlabel{fig:t6-gauss}{{6.1}{34}{El cuadrado $S_R$ contiene al disco $D_R$ y está contenido en $D_{R\sqrt 2}$. Las integrales de $e^{-(x^2+y^2)}$ sobre los discos se calculan en polares y encierran a $I(R)^2$}{figure.6.1}{}}
\newlabel{fig:t6-gauss@cref}{{[figure][1][6]6.1}{[1][33][]34}}
\newlabel{eje:6-simetria}{{6.1}{35}{}{ejercicio.6.1}{}}
\newlabel{eje:6-simetria@cref}{{[ejercicio][1][6]6.1}{[1][34][]35}}
\newlabel{eje:6-iteradas}{{6.2}{35}{}{ejercicio.6.2}{}}
\newlabel{eje:6-iteradas@cref}{{[ejercicio][2][6]6.2}{[1][35][]35}}
\newlabel{eje:6-rectangulos}{{6.3}{35}{}{ejercicio.6.3}{}}
\newlabel{eje:6-rectangulos@cref}{{[ejercicio][3][6]6.3}{[1][35][]35}}
\newlabel{eje:6-planteo}{{6.4}{35}{}{ejercicio.6.4}{}}
\newlabel{eje:6-planteo@cref}{{[ejercicio][4][6]6.4}{[1][35][]35}}
\newlabel{eje:6-region}{{6.5}{35}{}{ejercicio.6.5}{}}
\newlabel{eje:6-region@cref}{{[ejercicio][5][6]6.5}{[1][35][]35}}
\newlabel{eje:6-volumen-xy}{{6.6}{35}{}{ejercicio.6.6}{}}
\newlabel{eje:6-volumen-xy@cref}{{[ejercicio][6][6]6.6}{[1][35][]35}}
\newlabel{eje:6-orden-1}{{6.7}{35}{}{ejercicio.6.7}{}}
\newlabel{eje:6-orden-1@cref}{{[ejercicio][7][6]6.7}{[1][35][]35}}
\newlabel{eje:6-orden-2}{{6.8}{35}{}{ejercicio.6.8}{}}
\newlabel{eje:6-orden-2@cref}{{[ejercicio][8][6]6.8}{[1][35][]35}}
\newlabel{eje:6-polares}{{6.9}{35}{}{ejercicio.6.9}{}}
\newlabel{eje:6-polares@cref}{{[ejercicio][9][6]6.9}{[1][35][]35}}
\newlabel{eje:6-semiesfera}{{6.10}{35}{}{ejercicio.6.10}{}}
\newlabel{eje:6-semiesfera@cref}{{[ejercicio][10][6]6.10}{[1][35][]35}}
\newlabel{eje:6-centroide}{{6.11}{36}{}{ejercicio.6.11}{}}
\newlabel{eje:6-centroide@cref}{{[ejercicio][11][6]6.11}{[1][35][]36}}
\newlabel{eje:6-gauss-polar}{{6.12}{36}{}{ejercicio.6.12}{}}
\newlabel{eje:6-gauss-polar@cref}{{[ejercicio][12][6]6.12}{[1][36][]36}}
\newlabel{eje:6-capas}{{6.13}{36}{}{ejercicio.6.13}{}}
\newlabel{eje:6-capas@cref}{{[ejercicio][13][6]6.13}{[1][36][]36}}
\newlabel{eje:6-cota}{{6.14}{36}{}{ejercicio.6.14}{}}
\newlabel{eje:6-cota@cref}{{[ejercicio][14][6]6.14}{[1][36][]36}}
\newlabel{eje:6-gaussiana-normal}{{6.15}{36}{}{ejercicio.6.15}{}}
\newlabel{eje:6-gaussiana-normal@cref}{{[ejercicio][15][6]6.15}{[1][36][]36}}
\newlabel{ap:soluciones}{{A}{37}{Soluciones de los ejercicios}{appendix.A}{}}
\newlabel{ap:soluciones@cref}{{[appendix][1][2147483647]A}{[1][37][]37}}
\newlabel{sol:eje:1-vf}{{A.1}{37}{Tema 1}{section*.3}{}}
\newlabel{sol:eje:1-vf@cref}{{[subappendix][1][2147483647,1]A.1}{[1][37][]37}}
\newlabel{sol:eje:1-cosenos}{{A.1}{37}{Tema 1}{section*.4}{}}
\newlabel{sol:eje:1-cosenos@cref}{{[subappendix][1][2147483647,1]A.1}{[1][37][]37}}
\newlabel{sol:eje:1-operaciones}{{A.1}{38}{Tema 1}{section*.5}{}}
\newlabel{sol:eje:1-operaciones@cref}{{[subappendix][1][2147483647,1]A.1}{[1][38][]38}}
\newlabel{sol:eje:1-unitarios}{{A.1}{38}{Tema 1}{section*.6}{}}
\newlabel{sol:eje:1-unitarios@cref}{{[subappendix][1][2147483647,1]A.1}{[1][38][]38}}
\newlabel{sol:eje:1-angulos}{{A.1}{38}{Tema 1}{section*.7}{}}
\newlabel{sol:eje:1-angulos@cref}{{[subappendix][1][2147483647,1]A.1}{[1][38][]38}}
\newlabel{sol:eje:1-ortogonal-c}{{A.1}{39}{Tema 1}{section*.8}{}}
\newlabel{sol:eje:1-ortogonal-c@cref}{{[subappendix][1][2147483647,1]A.1}{[1][38][]39}}
\newlabel{sol:eje:1-proyeccion}{{A.1}{39}{Tema 1}{section*.9}{}}
\newlabel{sol:eje:1-proyeccion@cref}{{[subappendix][1][2147483647,1]A.1}{[1][39][]39}}
\newlabel{sol:eje:1-cubo}{{A.1}{39}{Tema 1}{section*.10}{}}
\newlabel{sol:eje:1-cubo@cref}{{[subappendix][1][2147483647,1]A.1}{[1][39][]39}}
\newlabel{sol:eje:1-vectorial}{{A.1}{39}{Tema 1}{section*.11}{}}
\newlabel{sol:eje:1-vectorial@cref}{{[subappendix][1][2147483647,1]A.1}{[1][39][]39}}
\newlabel{sol:eje:1-tetraedro}{{A.1}{40}{Tema 1}{section*.12}{}}
\newlabel{sol:eje:1-tetraedro@cref}{{[subappendix][1][2147483647,1]A.1}{[1][40][]40}}
\newlabel{sol:eje:1-polarizacion}{{A.1}{40}{Tema 1}{section*.13}{}}
\newlabel{sol:eje:1-polarizacion@cref}{{[subappendix][1][2147483647,1]A.1}{[1][40][]40}}
\newlabel{sol:eje:1-cs-cuadratica}{{A.1}{41}{Tema 1}{section*.14}{}}
\newlabel{sol:eje:1-cs-cuadratica@cref}{{[subappendix][1][2147483647,1]A.1}{[1][40][]41}}
\newlabel{sol:eje:1-rombo}{{A.1}{41}{Tema 1}{section*.15}{}}
\newlabel{sol:eje:1-rombo@cref}{{[subappendix][1][2147483647,1]A.1}{[1][41][]41}}
\newlabel{sol:eje:1-semicirculo}{{A.1}{41}{Tema 1}{section*.16}{}}
\newlabel{sol:eje:1-semicirculo@cref}{{[subappendix][1][2147483647,1]A.1}{[1][41][]41}}
\newlabel{sol:eje:1-lagrange-id}{{A.1}{42}{Tema 1}{section*.17}{}}
\newlabel{sol:eje:1-lagrange-id@cref}{{[subappendix][1][2147483647,1]A.1}{[1][42][]42}}
\newlabel{sol:eje:2-posicion}{{A.2}{42}{Tema 2}{section*.18}{}}
\newlabel{sol:eje:2-posicion@cref}{{[subappendix][2][2147483647,1]A.2}{[1][42][]42}}
\newlabel{sol:eje:2-planos}{{A.2}{42}{Tema 2}{section*.19}{}}
\newlabel{sol:eje:2-planos@cref}{{[subappendix][2][2147483647,1]A.2}{[1][42][]42}}
\newlabel{sol:eje:2-diedro}{{A.2}{43}{Tema 2}{section*.20}{}}
\newlabel{sol:eje:2-diedro@cref}{{[subappendix][2][2147483647,1]A.2}{[1][43][]43}}
\newlabel{sol:eje:2-interseccion}{{A.2}{43}{Tema 2}{section*.21}{}}
\newlabel{sol:eje:2-interseccion@cref}{{[subappendix][2][2147483647,1]A.2}{[1][43][]43}}
\newlabel{sol:eje:2-gauss}{{A.2}{43}{Tema 2}{section*.22}{}}
\newlabel{sol:eje:2-gauss@cref}{{[subappendix][2][2147483647,1]A.2}{[1][43][]43}}
\newlabel{sol:eje:2-gauss-casos}{{A.2}{44}{Tema 2}{section*.23}{}}
\newlabel{sol:eje:2-gauss-casos@cref}{{[subappendix][2][2147483647,1]A.2}{[1][44][]44}}
\newlabel{sol:eje:2-parametro}{{A.2}{44}{Tema 2}{section*.24}{}}
\newlabel{sol:eje:2-parametro@cref}{{[subappendix][2][2147483647,1]A.2}{[1][44][]44}}
\newlabel{sol:eje:2-areas}{{A.2}{44}{Tema 2}{section*.25}{}}
\newlabel{sol:eje:2-areas@cref}{{[subappendix][2][2147483647,1]A.2}{[1][44][]44}}
\newlabel{sol:eje:2-inversa}{{A.2}{45}{Tema 2}{section*.26}{}}
\newlabel{sol:eje:2-inversa@cref}{{[subappendix][2][2147483647,1]A.2}{[1][45][]45}}
\newlabel{sol:eje:2-produccion}{{A.2}{45}{Tema 2}{section*.27}{}}
\newlabel{sol:eje:2-produccion@cref}{{[subappendix][2][2147483647,1]A.2}{[1][45][]45}}
\newlabel{sol:eje:2-cramer}{{A.2}{45}{Tema 2}{section*.28}{}}
\newlabel{sol:eje:2-cramer@cref}{{[subappendix][2][2147483647,1]A.2}{[1][45][]45}}
\newlabel{sol:eje:2-distancia}{{A.2}{46}{Tema 2}{section*.29}{}}
\newlabel{sol:eje:2-distancia@cref}{{[subappendix][2][2147483647,1]A.2}{[1][46][]46}}
\newlabel{sol:eje:2-vandermonde}{{A.2}{46}{Tema 2}{section*.30}{}}
\newlabel{sol:eje:2-vandermonde@cref}{{[subappendix][2][2147483647,1]A.2}{[1][46][]46}}
\newlabel{sol:eje:2-compatibilidad}{{A.2}{47}{Tema 2}{section*.31}{}}
\newlabel{sol:eje:2-compatibilidad@cref}{{[subappendix][2][2147483647,1]A.2}{[1][47][]47}}
\newlabel{sol:eje:2-producto-det}{{A.2}{47}{Tema 2}{section*.32}{}}
\newlabel{sol:eje:2-producto-det@cref}{{[subappendix][2][2147483647,1]A.2}{[1][47][]47}}
\newlabel{sol:eje:3-conceptos}{{A.3}{48}{Tema 3}{section*.33}{}}
\newlabel{sol:eje:3-conceptos@cref}{{[subappendix][3][2147483647,1]A.3}{[1][48][]48}}
\newlabel{sol:eje:3-dominios}{{A.3}{48}{Tema 3}{section*.34}{}}
\newlabel{sol:eje:3-dominios@cref}{{[subappendix][3][2147483647,1]A.3}{[1][48][]48}}
\newlabel{sol:eje:3-niveles}{{A.3}{48}{Tema 3}{section*.35}{}}
\newlabel{sol:eje:3-niveles@cref}{{[subappendix][3][2147483647,1]A.3}{[1][48][]48}}
\newlabel{sol:eje:3-niveles-2}{{A.3}{49}{Tema 3}{section*.36}{}}
\newlabel{sol:eje:3-niveles-2@cref}{{[subappendix][3][2147483647,1]A.3}{[1][49][]49}}
\newlabel{sol:eje:3-onda}{{A.3}{49}{Tema 3}{section*.37}{}}
\newlabel{sol:eje:3-onda@cref}{{[subappendix][3][2147483647,1]A.3}{[1][49][]49}}
\newlabel{sol:eje:3-interseccion-superficies}{{A.3}{50}{Tema 3}{section*.38}{}}
\newlabel{sol:eje:3-interseccion-superficies@cref}{{[subappendix][3][2147483647,1]A.3}{[1][50][]50}}
\newlabel{sol:eje:3-limites-directos}{{A.3}{50}{Tema 3}{section*.39}{}}
\newlabel{sol:eje:3-limites-directos@cref}{{[subappendix][3][2147483647,1]A.3}{[1][50][]50}}
\newlabel{sol:eje:3-no-existe}{{A.3}{50}{Tema 3}{section*.40}{}}
\newlabel{sol:eje:3-no-existe@cref}{{[subappendix][3][2147483647,1]A.3}{[1][50][]50}}
\newlabel{sol:eje:3-polares}{{A.3}{51}{Tema 3}{section*.41}{}}
\newlabel{sol:eje:3-polares@cref}{{[subappendix][3][2147483647,1]A.3}{[1][50][]51}}
\newlabel{sol:eje:3-continuidad}{{A.3}{51}{Tema 3}{section*.42}{}}
\newlabel{sol:eje:3-continuidad@cref}{{[subappendix][3][2147483647,1]A.3}{[1][51][]51}}
\newlabel{sol:eje:3-lse}{{A.3}{51}{Tema 3}{section*.43}{}}
\newlabel{sol:eje:3-lse@cref}{{[subappendix][3][2147483647,1]A.3}{[1][51][]51}}
\newlabel{sol:eje:3-epsilon-delta}{{A.3}{52}{Tema 3}{section*.44}{}}
\newlabel{sol:eje:3-epsilon-delta@cref}{{[subappendix][3][2147483647,1]A.3}{[1][52][]52}}
\newlabel{sol:eje:3-composicion}{{A.3}{52}{Tema 3}{section*.45}{}}
\newlabel{sol:eje:3-composicion@cref}{{[subappendix][3][2147483647,1]A.3}{[1][52][]52}}
\newlabel{sol:eje:3-cuadratica}{{A.3}{52}{Tema 3}{section*.46}{}}
\newlabel{sol:eje:3-cuadratica@cref}{{[subappendix][3][2147483647,1]A.3}{[1][52][]52}}
\newlabel{sol:eje:3-trayectorias}{{A.3}{53}{Tema 3}{section*.47}{}}
\newlabel{sol:eje:3-trayectorias@cref}{{[subappendix][3][2147483647,1]A.3}{[1][53][]53}}
\newlabel{sol:eje:4-conceptos}{{A.4}{53}{Tema 4}{section*.48}{}}
\newlabel{sol:eje:4-conceptos@cref}{{[subappendix][4][2147483647,1]A.4}{[1][53][]53}}
\newlabel{sol:eje:4-parciales}{{A.4}{53}{Tema 4}{section*.49}{}}
\newlabel{sol:eje:4-parciales@cref}{{[subappendix][4][2147483647,1]A.4}{[1][53][]53}}
\newlabel{sol:eje:4-schwarz}{{A.4}{54}{Tema 4}{section*.50}{}}
\newlabel{sol:eje:4-schwarz@cref}{{[subappendix][4][2147483647,1]A.4}{[1][54][]54}}
\newlabel{sol:eje:4-tangente}{{A.4}{54}{Tema 4}{section*.51}{}}
\newlabel{sol:eje:4-tangente@cref}{{[subappendix][4][2147483647,1]A.4}{[1][54][]54}}
\newlabel{sol:eje:4-aproximacion}{{A.4}{54}{Tema 4}{section*.52}{}}
\newlabel{sol:eje:4-aproximacion@cref}{{[subappendix][4][2147483647,1]A.4}{[1][54][]54}}
\newlabel{sol:eje:4-gradiente}{{A.4}{55}{Tema 4}{section*.53}{}}
\newlabel{sol:eje:4-gradiente@cref}{{[subappendix][4][2147483647,1]A.4}{[1][55][]55}}
\newlabel{sol:eje:4-max-direccional}{{A.4}{55}{Tema 4}{section*.54}{}}
\newlabel{sol:eje:4-max-direccional@cref}{{[subappendix][4][2147483647,1]A.4}{[1][55][]55}}
\newlabel{sol:eje:4-tangente-nivel}{{A.4}{55}{Tema 4}{section*.55}{}}
\newlabel{sol:eje:4-tangente-nivel@cref}{{[subappendix][4][2147483647,1]A.4}{[1][55][]55}}
\newlabel{sol:eje:4-cadena-t}{{A.4}{56}{Tema 4}{section*.56}{}}
\newlabel{sol:eje:4-cadena-t@cref}{{[subappendix][4][2147483647,1]A.4}{[1][55][]56}}
\newlabel{sol:eje:4-cadena-uv}{{A.4}{56}{Tema 4}{section*.57}{}}
\newlabel{sol:eje:4-cadena-uv@cref}{{[subappendix][4][2147483647,1]A.4}{[1][56][]56}}
\newlabel{sol:eje:4-softmax}{{A.4}{56}{Tema 4}{section*.58}{}}
\newlabel{sol:eje:4-softmax@cref}{{[subappendix][4][2147483647,1]A.4}{[1][56][]56}}
\newlabel{sol:eje:4-polares}{{A.4}{57}{Tema 4}{section*.59}{}}
\newlabel{sol:eje:4-polares@cref}{{[subappendix][4][2147483647,1]A.4}{[1][56][]57}}
\newlabel{sol:eje:4-euler}{{A.4}{57}{Tema 4}{section*.60}{}}
\newlabel{sol:eje:4-euler@cref}{{[subappendix][4][2147483647,1]A.4}{[1][57][]57}}
\newlabel{sol:eje:4-material}{{A.4}{57}{Tema 4}{section*.61}{}}
\newlabel{sol:eje:4-material@cref}{{[subappendix][4][2147483647,1]A.4}{[1][57][]57}}
\newlabel{sol:eje:4-no-diferenciable}{{A.4}{58}{Tema 4}{section*.62}{}}
\newlabel{sol:eje:4-no-diferenciable@cref}{{[subappendix][4][2147483647,1]A.4}{[1][57][]58}}
\newlabel{sol:eje:5-conceptos}{{A.5}{58}{Tema 5}{section*.63}{}}
\newlabel{sol:eje:5-conceptos@cref}{{[subappendix][5][2147483647,1]A.5}{[1][58][]58}}
\newlabel{sol:eje:5-clasificar-1}{{A.5}{58}{Tema 5}{section*.64}{}}
\newlabel{sol:eje:5-clasificar-1@cref}{{[subappendix][5][2147483647,1]A.5}{[1][58][]58}}
\newlabel{sol:eje:5-clasificar-2}{{A.5}{59}{Tema 5}{section*.65}{}}
\newlabel{sol:eje:5-clasificar-2@cref}{{[subappendix][5][2147483647,1]A.5}{[1][59][]59}}
\newlabel{sol:eje:5-clasificar-3}{{A.5}{59}{Tema 5}{section*.66}{}}
\newlabel{sol:eje:5-clasificar-3@cref}{{[subappendix][5][2147483647,1]A.5}{[1][59][]59}}
\newlabel{sol:eje:5-openstax}{{A.5}{60}{Tema 5}{section*.67}{}}
\newlabel{sol:eje:5-openstax@cref}{{[subappendix][5][2147483647,1]A.5}{[1][59][]60}}
\newlabel{sol:eje:5-absolutos}{{A.5}{60}{Tema 5}{section*.68}{}}
\newlabel{sol:eje:5-absolutos@cref}{{[subappendix][5][2147483647,1]A.5}{[1][60][]60}}
\newlabel{sol:eje:5-disco}{{A.5}{60}{Tema 5}{section*.69}{}}
\newlabel{sol:eje:5-disco@cref}{{[subappendix][5][2147483647,1]A.5}{[1][60][]60}}
\newlabel{sol:eje:5-caja}{{A.5}{60}{Tema 5}{section*.70}{}}
\newlabel{sol:eje:5-caja@cref}{{[subappendix][5][2147483647,1]A.5}{[1][60][]60}}
\newlabel{sol:eje:5-hiperbola}{{A.5}{61}{Tema 5}{section*.71}{}}
\newlabel{sol:eje:5-hiperbola@cref}{{[subappendix][5][2147483647,1]A.5}{[1][61][]61}}
\newlabel{sol:eje:5-resistencias}{{A.5}{61}{Tema 5}{section*.72}{}}
\newlabel{sol:eje:5-resistencias@cref}{{[subappendix][5][2147483647,1]A.5}{[1][61][]61}}
\newlabel{sol:eje:5-minimos-cuadrados}{{A.5}{62}{Tema 5}{section*.73}{}}
\newlabel{sol:eje:5-minimos-cuadrados@cref}{{[subappendix][5][2147483647,1]A.5}{[1][61][]62}}
\newlabel{sol:eje:5-degenerado}{{A.5}{62}{Tema 5}{section*.74}{}}
\newlabel{sol:eje:5-degenerado@cref}{{[subappendix][5][2147483647,1]A.5}{[1][62][]62}}
\newlabel{sol:eje:5-mle}{{A.5}{62}{Tema 5}{section*.75}{}}
\newlabel{sol:eje:5-mle@cref}{{[subappendix][5][2147483647,1]A.5}{[1][62][]62}}
\newlabel{sol:eje:5-media-geometrica}{{A.5}{63}{Tema 5}{section*.76}{}}
\newlabel{sol:eje:5-media-geometrica@cref}{{[subappendix][5][2147483647,1]A.5}{[1][62][]63}}
\newlabel{sol:eje:5-cs-lagrange}{{A.5}{63}{Tema 5}{section*.77}{}}
\newlabel{sol:eje:5-cs-lagrange@cref}{{[subappendix][5][2147483647,1]A.5}{[1][63][]63}}
\newlabel{sol:eje:6-simetria}{{A.6}{63}{Tema 6}{section*.78}{}}
\newlabel{sol:eje:6-simetria@cref}{{[subappendix][6][2147483647,1]A.6}{[1][63][]63}}
\newlabel{sol:eje:6-iteradas}{{A.6}{64}{Tema 6}{section*.79}{}}
\newlabel{sol:eje:6-iteradas@cref}{{[subappendix][6][2147483647,1]A.6}{[1][64][]64}}
\newlabel{sol:eje:6-rectangulos}{{A.6}{64}{Tema 6}{section*.80}{}}
\newlabel{sol:eje:6-rectangulos@cref}{{[subappendix][6][2147483647,1]A.6}{[1][64][]64}}
\newlabel{sol:eje:6-planteo}{{A.6}{64}{Tema 6}{section*.81}{}}
\newlabel{sol:eje:6-planteo@cref}{{[subappendix][6][2147483647,1]A.6}{[1][64][]64}}
\newlabel{sol:eje:6-region}{{A.6}{65}{Tema 6}{section*.82}{}}
\newlabel{sol:eje:6-region@cref}{{[subappendix][6][2147483647,1]A.6}{[1][65][]65}}
\newlabel{sol:eje:6-volumen-xy}{{A.6}{65}{Tema 6}{section*.83}{}}
\newlabel{sol:eje:6-volumen-xy@cref}{{[subappendix][6][2147483647,1]A.6}{[1][65][]65}}
\newlabel{sol:eje:6-orden-1}{{A.6}{65}{Tema 6}{section*.84}{}}
\newlabel{sol:eje:6-orden-1@cref}{{[subappendix][6][2147483647,1]A.6}{[1][65][]65}}
\newlabel{sol:eje:6-orden-2}{{A.6}{66}{Tema 6}{section*.85}{}}
\newlabel{sol:eje:6-orden-2@cref}{{[subappendix][6][2147483647,1]A.6}{[1][66][]66}}
\newlabel{sol:eje:6-polares}{{A.6}{66}{Tema 6}{section*.86}{}}
\newlabel{sol:eje:6-polares@cref}{{[subappendix][6][2147483647,1]A.6}{[1][66][]66}}
\newlabel{sol:eje:6-semiesfera}{{A.6}{66}{Tema 6}{section*.87}{}}
\newlabel{sol:eje:6-semiesfera@cref}{{[subappendix][6][2147483647,1]A.6}{[1][66][]66}}
\newlabel{sol:eje:6-centroide}{{A.6}{66}{Tema 6}{section*.88}{}}
\newlabel{sol:eje:6-centroide@cref}{{[subappendix][6][2147483647,1]A.6}{[1][66][]66}}
\newlabel{sol:eje:6-gauss-polar}{{A.6}{67}{Tema 6}{section*.89}{}}
\newlabel{sol:eje:6-gauss-polar@cref}{{[subappendix][6][2147483647,1]A.6}{[1][67][]67}}
\newlabel{sol:eje:6-capas}{{A.6}{67}{Tema 6}{section*.90}{}}
\newlabel{sol:eje:6-capas@cref}{{[subappendix][6][2147483647,1]A.6}{[1][67][]67}}
\newlabel{sol:eje:6-cota}{{A.6}{67}{Tema 6}{section*.91}{}}
\newlabel{sol:eje:6-cota@cref}{{[subappendix][6][2147483647,1]A.6}{[1][67][]67}}
\newlabel{sol:eje:6-gaussiana-normal}{{A.6}{68}{Tema 6}{section*.92}{}}
\newlabel{sol:eje:6-gaussiana-normal@cref}{{[subappendix][6][2147483647,1]A.6}{[1][68][]68}}
\makeatother

\begin{document}
\frontmatter
\begin{titlepage}
\centering
\vspace*{2.5cm}
{\large\scshape Anatomía de un LLM}\par
\vspace{0.4cm}
{\normalsize Fundamentos matemáticos de los modelos de lenguaje}\par
\vspace{2.2cm}
\rule{\textwidth}{0.8pt}\par\vspace{0.5cm}
{\small\scshape Fase 0 \textperiodcentered\ Preliminares matemáticos}\par\vspace{0.3cm}
{\Huge\bfseries Módulo 0.7\par}\vspace{0.3cm}
{\LARGE Vectores y funciones de varias variables\par}
\vspace{0.5cm}\rule{\textwidth}{0.4pt}\par
\vspace{1.6cm}
{\large Resumen del módulo}\par
\vfill
{\small Edición de octubre de 2026}\par
\end{titlepage}

\chapter*{Sobre este resumen}
\addcontentsline{toc}{chapter}{Sobre este resumen}
\markboth{Sobre este resumen}{}

Este documento reúne lo indispensable de los apuntes del Módulo 0.7: el texto expositivo, todas las definiciones y todos los resultados (teoremas, proposiciones, lemas y corolarios), las tablas y las observaciones de contenido. Se han suprimido las demostraciones, los ejemplos, las lecturas, las figuras, los ejercicios y sus soluciones. La numeración de temas, secciones, enunciados y ecuaciones es la misma que en los apuntes, de modo que cualquier referencia puede buscarse allí; las que remiten a elementos suprimidos (un ejemplo, una figura, una ecuación de una demostración) conservan su número en los apuntes.

\section*{Notación}
\begin{center}
\begin{tabular}{@{}ll@{}}
\toprule
Símbolo & Significado \\
\midrule
$\vu=(u_1,u_2,u_3)$, $\vi,\vj,\vk$ & vector de $\R^3$ y base canónica \\
$\vu\cdot\vv$, $\norm{\vu}$ & producto escalar y norma euclídea \\
$\vu\times\vv$ & producto vectorial \\
$\proj_{\vv}\vu$ & proyección ortogonal de $\vu$ sobre $\vv$ \\
$\det A$, $A^{-1}$ & determinante e inversa de una matriz cuadrada \\
$f_x=\pd fx$, $f_{xy}=\frac{\partial^2 f}{\partial y\,\partial x}$ & derivadas parciales primera y segunda \\
$\nabla f$, $D_{\vu}f$ & gradiente y derivada direccional según el vector unitario $\vu$ \\
$H_f$ & matriz hessiana de $f$ \\
$\iint_R f\dA$ & integral doble sobre la región $R$ \\
$(r,\theta)$ & coordenadas polares: $x=r\cos\theta$, $y=r\sin\theta$ \\
\bottomrule
\end{tabular}
\end{center}


\tableofcontents

\mainmatter
% =====================================================================
\chapter{Vectores en el plano y en el espacio}
\label{tema:1}
% =====================================================================

Un vector es, a la vez, una lista de números y una flecha con longitud y dirección. Este tema relaciona las dos caras: las operaciones se definen con coordenadas y se interpretan geométricamente. El producto escalar permite medir longitudes y ángulos; el producto vectorial, áreas y volúmenes. La presentación sigue las sesiones~1 a~7 de la unidad~1 del curso 18.02SC del MIT~\cite{mit1802}, las secciones 2.1--2.4 de OpenStax~\cite{openstax3} y la sección~11.1 de Strang~\cite{strangcalc}.

\section{Vectores y operaciones}

\begin{definicion}[Vectores de $\R^n$]\label{def:vector}
Para $n=2$ o $n=3$, un \emph{vector} de $\R^n$ es una lista ordenada $\vu=(u_1,\dots,u_n)$ de números reales, sus \emph{componentes}. Para $\vu,\vv\in\R^n$ y $\lambda\in\R$ se definen la \emph{suma} y el \emph{producto por un escalar} componente a componente:
\[
\vu+\vv=(u_1+v_1,\dots,u_n+v_n),\qquad \lambda\vu=(\lambda u_1,\dots,\lambda u_n).
\]
El vector nulo es $\vcero=(0,\dots,0)$ y $-\vu=(-1)\vu$. En $\R^3$ se escribe $\vi=(1,0,0)$, $\vj=(0,1,0)$, $\vk=(0,0,1)$, de modo que $(u_1,u_2,u_3)=u_1\vi+u_2\vj+u_3\vk$.
\end{definicion}

Las propiedades de la suma y del producto por escalares (asociativa, conmutativa, distributivas, neutro y opuesto) se heredan de las de los números reales, porque se opera componente a componente. Un \emph{punto} y un \emph{vector} se escriben igual; la diferencia es de interpretación. Si $P$ y $Q$ son puntos, el vector $\overrightarrow{PQ}=Q-P$ es la flecha que va de $P$ a $Q$; dos flechas con la misma longitud, dirección y sentido representan el mismo vector, aunque partan de puntos distintos.

Geométricamente, $\vu+\vv$ es la diagonal del paralelogramo de lados $\vu$ y $\vv$ (o, equivalentemente, la flecha que se obtiene colocando $\vv$ a continuación de $\vu$), y $\lambda\vu$ es una flecha paralela a $\vu$, $\abs{\lambda}$ veces más larga y con el sentido invertido si $\lambda<0$.

\addtocounter{definicion}{1}%

\section{Producto escalar, norma y ángulo}

\begin{definicion}[Producto escalar y norma]\label{def:escalar}
El \emph{producto escalar} de $\vu,\vv\in\R^n$ es el número
\[
\vu\cdot\vv=u_1v_1+\dots+u_nv_n=\sum_{i=1}^{n}u_iv_i .
\]
La \emph{norma} (o longitud) de $\vu$ es $\norm{\vu}=\sqrt{\vu\cdot\vu}=\sqrt{u_1^2+\dots+u_n^2}$, y la \emph{distancia} entre $\vu$ y $\vv$ es $\norm{\vu-\vv}$. Un vector de norma~1 es \emph{unitario}; si $\vu\neq\vcero$, el vector $\vu/\norm{\vu}$ es el unitario en la dirección de $\vu$.
\end{definicion}

En $\R^2$ y $\R^3$ la norma es la longitud de la flecha, por el teorema de Pitágoras (aplicado dos veces en $\R^3$).

\begin{proposicion}[Propiedades del producto escalar]\label{prop:escalar}
Para todos $\vu,\vv,\vw\in\R^n$ y $\lambda\in\R$:
\begin{enumerate}
  \item simetría: $\vu\cdot\vv=\vv\cdot\vu$;
  \item linealidad: $(\vu+\vw)\cdot\vv=\vu\cdot\vv+\vw\cdot\vv$ y $(\lambda\vu)\cdot\vv=\lambda(\vu\cdot\vv)$;
  \item positividad: $\vu\cdot\vu\ge 0$, con igualdad si y solo si $\vu=\vcero$;
  \item $\norm{\lambda\vu}=\abs{\lambda}\,\norm{\vu}$.
\end{enumerate}
\end{proposicion}

Desarrollando con (i) y (ii) se obtiene la identidad, que se usará muchas veces,
\begin{equation}\label{eq:t1-desarrollo}
\norm{\vu\pm\vv}^2=(\vu\pm\vv)\cdot(\vu\pm\vv)=\norm{\vu}^2\pm2\,\vu\cdot\vv+\norm{\vv}^2 .
\end{equation}

\begin{teorema}[Desigualdad de Cauchy--Schwarz]\label{teo:cs}
Para todos $\vu,\vv\in\R^n$, $\abs{\vu\cdot\vv}\le\norm{\vu}\,\norm{\vv}$. La igualdad se da si y solo si uno de los vectores es múltiplo del otro.
\end{teorema}

\begin{corolario}[Desigualdad triangular]\label{cor:triangular}
Para todos $\vu,\vv\in\R^n$, $\norm{\vu+\vv}\le\norm{\vu}+\norm{\vv}$.
\end{corolario}

Por Cauchy--Schwarz, si $\vu,\vv\neq\vcero$ el cociente $\vu\cdot\vv/(\norm{\vu}\norm{\vv})$ está en $[-1,1]$, lo que justifica la definición siguiente.

\begin{definicion}[Ángulo y ortogonalidad]\label{def:angulo}
El \emph{ángulo} entre dos vectores no nulos $\vu$ y $\vv$ es el único $\theta\in[0,\pi]$ tal que
\[
\cos\theta=\frac{\vu\cdot\vv}{\norm{\vu}\,\norm{\vv}} .
\]
Dos vectores son \emph{ortogonales} (o perpendiculares), $\vu\perp\vv$, si $\vu\cdot\vv=0$.
\end{definicion}

\begin{proposicion}[El ángulo algebraico es el geométrico]\label{prop:coseno}
Si $\vu,\vv\neq\vcero$ forman, como flechas con origen común, un ángulo geométrico $\theta$, entonces $\vu\cdot\vv=\norm{\vu}\norm{\vv}\cos\theta$.
\end{proposicion}

Así, el signo de $\vu\cdot\vv$ informa del ángulo: positivo si es agudo, nulo si es recto y negativo si es obtuso. Si $\vu\perp\vv$, la identidad~\eqref{eq:t1-desarrollo} se reduce al \emph{teorema de Pitágoras} $\norm{\vu+\vv}^2=\norm{\vu}^2+\norm{\vv}^2$.

\addtocounter{definicion}{1}%

\section{Proyecciones}

\begin{definicion}[Proyección ortogonal]\label{def:proyeccion}
Sean $\vu,\vv\in\R^n$ con $\vv\neq\vcero$. La \emph{componente escalar} de $\vu$ en la dirección de $\vv$ es $\operatorname{comp}_{\vv}\vu=\vu\cdot\vv/\norm{\vv}$, y la \emph{proyección ortogonal} de $\vu$ sobre $\vv$ es el vector
\[
\proj_{\vv}\vu=\frac{\vu\cdot\vv}{\vv\cdot\vv}\,\vv .
\]
\end{definicion}

\begin{proposicion}[Descomposición ortogonal]\label{prop:descomposicion}
Sea $\vv\neq\vcero$. Todo $\vu$ se escribe de forma única como $\vu=\vu_\parallel+\vu_\perp$ con $\vu_\parallel$ múltiplo de $\vv$ y $\vu_\perp\perp\vv$; además, $\vu_\parallel=\proj_{\vv}\vu$. Entre todos los múltiplos $t\vv$, el más cercano a $\vu$ es $\proj_{\vv}\vu$.
\end{proposicion}

\addtocounter{figure}{1}%

\addtocounter{definicion}{1}%

\section{Producto vectorial y producto mixto}

\begin{definicion}[Producto vectorial]\label{def:vectorial}
El \emph{producto vectorial} de $\vu,\vv\in\R^3$ es el vector
\[
\vu\times\vv=(u_2v_3-u_3v_2,\ u_3v_1-u_1v_3,\ u_1v_2-u_2v_1)
=\begin{vmatrix}\vi&\vj&\vk\\u_1&u_2&u_3\\v_1&v_2&v_3\end{vmatrix},
\]
donde el determinante es simbólico y se desarrolla por la primera fila (\cref{def:det3}).
\end{definicion}

\begin{proposicion}[Propiedades del producto vectorial]\label{prop:vectorial}
Para todos $\vu,\vv,\vw\in\R^3$ y $\lambda\in\R$:
\begin{enumerate}
  \item anticonmutatividad: $\vv\times\vu=-\vu\times\vv$; en particular, $\vu\times\vu=\vcero$;
  \item linealidad: $(\vu+\vw)\times\vv=\vu\times\vv+\vw\times\vv$ y $(\lambda\vu)\times\vv=\lambda(\vu\times\vv)$;
  \item ortogonalidad: $(\vu\times\vv)\cdot\vu=(\vu\times\vv)\cdot\vv=0$;
  \item $\norm{\vu\times\vv}=\norm{\vu}\norm{\vv}\sin\theta$, donde $\theta$ es el ángulo entre $\vu$ y $\vv$; por tanto, $\norm{\vu\times\vv}$ es el área del paralelogramo de lados $\vu$ y $\vv$;
  \item $\vi\times\vj=\vk$, $\vj\times\vk=\vi$, $\vk\times\vi=\vj$.
\end{enumerate}
\end{proposicion}

La dirección de $\vu\times\vv$ es, pues, perpendicular al plano que contienen $\vu$ y $\vv$, y su sentido viene dado por la \emph{regla de la mano derecha}: $(\vu,\vv,\vu\times\vv)$ está orientado como $(\vi,\vj,\vk)$. El producto vectorial no es asociativo: $(\vi\times\vi)\times\vj=\vcero$, mientras que $\vi\times(\vi\times\vj)=\vi\times\vk=-\vj$.

\begin{definicion}[Producto mixto]\label{def:mixto}
El \emph{producto mixto} de $\vu,\vv,\vw\in\R^3$ es el número $\vu\cdot(\vv\times\vw)$.
\end{definicion}

\begin{proposicion}[Volumen]\label{prop:volumen}
$\abs{\vu\cdot(\vv\times\vw)}$ es el volumen del paralelepípedo de aristas $\vu$, $\vv$ y $\vw$ con un vértice común. En particular, el producto mixto es nulo si y solo si los tres vectores están en un mismo plano por el origen.
\end{proposicion}

\addtocounter{definicion}{1}%

\addtocounter{definicion}{1}%

\addtocounter{definicion}{1}%


% =====================================================================
\chapter{Rectas, planos, sistemas lineales y determinantes}
\label{tema:2}
% =====================================================================

Una ecuación lineal en tres incógnitas describe un plano, y resolver un sistema de tres ecuaciones es cortar tres planos. Este tema escribe rectas y planos con vectores, resuelve sistemas pequeños por eliminación gaussiana e interpreta el determinante como un área o un volumen con signo. Todo se limita a tamaños $2\times2$ y $3\times3$: la teoría general (rango, independencia lineal, determinantes de orden $n$) es materia de los módulos~2 y~3 del curso. La presentación sigue las sesiones~5--16 de la unidad~1 de 18.02SC~\cite{mit1802}, la §2.5 de OpenStax~\cite{openstax3}, las secciones~11.2--11.4 de Strang~\cite{strangcalc} y, para la eliminación, el capítulo~2 de su \emph{Introduction to Linear Algebra}~\cite{strangla}.

\section{Rectas y planos}

\begin{definicion}[Recta]\label{def:recta}
La \emph{recta} que pasa por el punto $P_0$ con vector director $\vv\neq\vcero$ es el conjunto de puntos
\[
\vr(t)=P_0+t\,\vv,\qquad t\in\R .
\]
En coordenadas, $x=x_0+tv_1$, $y=y_0+tv_2$, $z=z_0+tv_3$ (\emph{ecuaciones paramétricas}). La recta que pasa por $P$ y $Q$ tiene vector director $\overrightarrow{PQ}$.
\end{definicion}

En el plano $\R^2$ una recta admite además una ecuación \emph{implícita} $ax+by=c$; el vector $(a,b)$ es perpendicular a la recta, porque si $(x_1,y_1)$ y $(x_2,y_2)$ son dos de sus puntos, restando las ecuaciones queda $(a,b)\cdot(x_2-x_1,y_2-y_1)=0$. En el espacio la ecuación implícita de una sola ecuación describe un plano.

\begin{definicion}[Plano]\label{def:plano}
El \emph{plano} que pasa por $P_0=(x_0,y_0,z_0)$ con \emph{vector normal} $\vn=(a,b,c)\neq\vcero$ es el conjunto de los puntos $X=(x,y,z)$ tales que $\vn\cdot(X-P_0)=0$, es decir,
\[
a(x-x_0)+b(y-y_0)+c(z-z_0)=0\quad\Longleftrightarrow\quad ax+by+cz=d,\qquad d=\vn\cdot P_0 .
\]
\end{definicion}

Recíprocamente, toda ecuación $ax+by+cz=d$ con $(a,b,c)\neq\vcero$ es un plano de normal $(a,b,c)$: basta tomar uno de sus puntos como $P_0$. Para el plano que pasa por tres puntos no alineados $P$, $Q$, $R$ se toma como normal $\vn=\overrightarrow{PQ}\times\overrightarrow{PR}$, que es perpendicular a ambos vectores (\cref{prop:vectorial}).

\addtocounter{definicion}{1}%

\begin{proposicion}[Posición relativa de dos planos]\label{prop:dos-planos}
Sean $\Pi_1:\vn_1\cdot X=d_1$ y $\Pi_2:\vn_2\cdot X=d_2$.
\begin{enumerate}
  \item Si $\vn_1$ y $\vn_2$ son proporcionales, los planos son paralelos: coinciden si las ecuaciones son proporcionales y no tienen puntos comunes en caso contrario.
  \item Si no lo son, se cortan en una recta de vector director $\vn_1\times\vn_2$.
\end{enumerate}
El \emph{ángulo} entre $\Pi_1$ y $\Pi_2$ es el ángulo agudo (o recto) $\varphi$ entre sus normales: $\cos\varphi=\abs{\vn_1\cdot\vn_2}/(\norm{\vn_1}\norm{\vn_2})$.
\end{proposicion}

\addtocounter{definicion}{1}%

Para cortar una recta $\vr(t)=P_0+t\vv$ con un plano $\vn\cdot X=d$ se sustituye $\vr(t)$ en la ecuación del plano y se despeja $t$. Por ejemplo, la recta $(1+t,\,t,\,t)$ y el plano $x+2y+3z=7$ dan $1+6t=7$, $t=1$, y el punto de corte es $(2,1,1)$. Si $\vn\cdot\vv=0$, la recta es paralela al plano: o no lo corta o está contenida en él.

\section{Sistemas lineales y eliminación gaussiana}\label{sec:t2-gauss}

\begin{definicion}[Sistema lineal y matriz]\label{def:sistema}
Un sistema de $m$ ecuaciones lineales con $n$ incógnitas ($m,n\le3$ en este tema) se escribe $A\vx=\vb$, donde $A=(a_{ij})$ es la \emph{matriz de coeficientes}, $\vx$ el vector de incógnitas y $\vb$ el de términos independientes. El producto $A\vx$ es el vector cuya componente $i$ es el producto escalar de la fila $i$ de $A$ por $\vx$. El sistema es \emph{compatible determinado} si tiene una única solución, \emph{compatible indeterminado} si tiene infinitas e \emph{incompatible} si no tiene ninguna.
\end{definicion}

Más en general, el producto de matrices $AB$ (definido cuando el número de columnas de $A$ coincide con el de filas de $B$) tiene como elemento $(i,j)$ el producto escalar de la fila $i$ de $A$ por la columna $j$ de $B$. La matriz identidad $I$ tiene unos en la diagonal y ceros fuera, y cumple $AI=IA=A$.

\begin{definicion}[Operaciones elementales]\label{def:elementales}
Las \emph{operaciones elementales} sobre las ecuaciones de un sistema son: (E1) intercambiar dos ecuaciones; (E2) multiplicar una ecuación por un número no nulo; (E3) sumar a una ecuación un múltiplo de otra.
\end{definicion}

\begin{proposicion}\label{prop:elementales}
Las operaciones elementales no cambian el conjunto de soluciones de un sistema.
\end{proposicion}

El \emph{método de Gauss} aplica operaciones elementales hasta llegar a un sistema \emph{escalonado} (cada ecuación tiene al menos una incógnita inicial menos que la anterior) y lo resuelve por \emph{sustitución regresiva}, de abajo arriba. En la práctica se trabaja con la \emph{matriz ampliada} $(A\mid\vb)$.

\addtocounter{definicion}{1}%

Si al escalonar aparece una fila $(0\ 0\ 0\mid c)$ con $c\neq0$, el sistema es incompatible; si aparece una fila entera de ceros y quedan menos ecuaciones que incógnitas, las incógnitas que no encabezan ninguna fila son \emph{libres} y el sistema es compatible indeterminado. Geométricamente, en un sistema $3\times3$ cada ecuación es un plano: la solución única corresponde a tres planos con un único punto común; la indeterminación, a planos que comparten una recta (o un plano); la incompatibilidad, a planos sin punto común a los tres.

\section{Determinantes de orden 2 y 3}

\begin{definicion}[Determinantes]\label{def:det3}
\[
\det\begin{pmatrix}a&b\\c&d\end{pmatrix}=\begin{vmatrix}a&b\\c&d\end{vmatrix}=ad-bc,
\qquad
\begin{vmatrix}a_1&a_2&a_3\\b_1&b_2&b_3\\c_1&c_2&c_3\end{vmatrix}
=a_1\begin{vmatrix}b_2&b_3\\c_2&c_3\end{vmatrix}-a_2\begin{vmatrix}b_1&b_3\\c_1&c_3\end{vmatrix}+a_3\begin{vmatrix}b_1&b_2\\c_1&c_2\end{vmatrix}.
\]
\end{definicion}

Comparando con la \cref{def:vectorial}, el determinante $3\times3$ de filas $\va,\vb,\mathbf c$ es exactamente el producto mixto $\va\cdot(\vb\times\mathbf c)$.

\begin{teorema}[Interpretación geométrica]\label{teo:det-area}
\begin{enumerate}
  \item $\abs{ad-bc}$ es el área del paralelogramo de lados $(a,b)$ y $(c,d)$ en $\R^2$.
  \item El valor absoluto del determinante $3\times3$ de filas $\va,\vb,\mathbf c$ es el volumen del paralelepípedo de aristas $\va,\vb,\mathbf c$.
\end{enumerate}
\end{teorema}

\begin{proposicion}[Propiedades de los determinantes]\label{prop:det-propiedades}
Para matrices $2\times2$ y $3\times3$:
\begin{enumerate}
  \item intercambiar dos filas cambia el signo del determinante; si dos filas son iguales, el determinante es nulo;
  \item multiplicar una fila por $\lambda$ multiplica el determinante por $\lambda$;
  \item sumar a una fila un múltiplo de otra no cambia el determinante;
  \item $\det(A^{\mathsf T})=\det A$ y $\det(AB)=\det A\cdot\det B$.
\end{enumerate}
\end{proposicion}

Para orden~2 todas las propiedades se comprueban con la fórmula (véase el \cref{eje:2-producto-det}); para orden~3 se deducen de las propiedades del producto mixto, y la demostración general, para cualquier orden, se da en el módulo~3 del curso. La propiedad~(iii) dice que la eliminación gaussiana con operaciones (E3) conserva el determinante: el de una matriz escalonada es el producto de los elementos de su diagonal. En el \cref{ej:t2-gauss}, $\det A=1\cdot1\cdot2=2$.

\begin{teorema}[Sistemas $2\times2$ y regla de Cramer]\label{teo:cramer}
El sistema $ax+by=e$, $cx+dy=f$ tiene solución única si y solo si $\Delta=ad-bc\neq0$, y en ese caso
\[
x=\frac{1}{\Delta}\begin{vmatrix}e&b\\f&d\end{vmatrix},\qquad y=\frac{1}{\Delta}\begin{vmatrix}a&e\\c&f\end{vmatrix}.
\]
\end{teorema}

Lo mismo vale en orden~3: un sistema $3\times3$ tiene solución única si y solo si $\det A\neq0$, es decir, si las tres normales (las filas de $A$) no son coplanarias (\cref{prop:volumen}); la regla de Cramer sustituye la columna $j$ de $A$ por $\vb$ para obtener la incógnita $j$.

\begin{definicion}[Matriz inversa]\label{def:inversa}
Una matriz cuadrada $A$ es \emph{invertible} si existe $A^{-1}$ con $AA^{-1}=A^{-1}A=I$. Entonces el sistema $A\vx=\vb$ tiene la única solución $\vx=A^{-1}\vb$.
\end{definicion}

Para orden~2, si $ad-bc\neq0$,
\begin{equation}\label{eq:t2-inversa}
\begin{pmatrix}a&b\\c&d\end{pmatrix}^{-1}=\frac{1}{ad-bc}\begin{pmatrix}d&-b\\-c&a\end{pmatrix},
\end{equation}
como se comprueba multiplicando. Una matriz es invertible si y solo si su determinante es no nulo.

\addtocounter{definicion}{1}%

\addtocounter{definicion}{1}%


% =====================================================================
\chapter{Funciones de varias variables, límites y continuidad}
\label{tema:3}
% =====================================================================

Una función de varias variables asigna un número a cada punto del plano o del espacio: la altitud de un terreno, la temperatura de una habitación o la pérdida de una red neuronal en función de sus parámetros. Este tema enseña a visualizarlas mediante gráficas y curvas de nivel y extiende a ellas las nociones de límite y continuidad. La novedad respecto a una variable es que un punto puede alcanzarse por infinitos caminos, y el límite debe ser el mismo por todos ellos. La presentación sigue las §4.1--4.2 de OpenStax~\cite{openstax3}, las sesiones~24--25 de 18.02SC~\cite{mit1802}, el capítulo~8 de Apostol~\cite{apostol2} y la §13.1 de Strang~\cite{strangcalc}.

\section{Funciones, dominio y gráfica}

\begin{definicion}[Función de varias variables]\label{def:funcion-n}
Una \emph{función real de $n$ variables} es una regla $f$ que asigna a cada punto $\vx=(x_1,\dots,x_n)$ de un conjunto $D\subseteq\R^n$, su \emph{dominio}, un único número real $f(\vx)$. Si no se indica, el dominio es el mayor conjunto en el que la fórmula tiene sentido. La \emph{imagen} es $\{f(\vx):\vx\in D\}$. Para $n=2$, la \emph{gráfica} de $f$ es la superficie $\{(x,y,z)\in\R^3: (x,y)\in D,\ z=f(x,y)\}$.
\end{definicion}

\addtocounter{definicion}{1}%

\section{Curvas y superficies de nivel}

\begin{definicion}[Conjuntos de nivel]\label{def:nivel}
Sea $f:D\subseteq\R^n\to\R$ y $c\in\R$. El \emph{conjunto de nivel} $c$ de $f$ es $\{\vx\in D: f(\vx)=c\}$. Para $n=2$ se llama \emph{curva de nivel}, y para $n=3$, \emph{superficie de nivel}. Un \emph{mapa de contorno} dibuja varias curvas de nivel, normalmente para valores de $c$ igualmente espaciados.
\end{definicion}

La curva de nivel $c$ es la proyección sobre el plano $xy$ del corte de la gráfica con el plano horizontal $z=c$, como las curvas de un mapa topográfico. Dos curvas de nivel distintas nunca se cortan, porque $f$ asigna un único valor a cada punto. Donde las curvas de nivel (equiespaciadas) están más juntas, la función cambia más deprisa.

\addtocounter{figure}{1}%

\addtocounter{definicion}{1}%

\section{Límites}

En $\R^n$ la cercanía se mide con la norma euclídea: la \emph{bola abierta} de centro $\va$ y radio $\delta>0$ es $B(\va,\delta)=\{\vx:\norm{\vx-\va}<\delta\}$. Se dice que $\va$ es un \emph{punto de acumulación} de $D$ si toda bola centrada en $\va$ contiene puntos de $D$ distintos de $\va$.

\begin{definicion}[Límite]\label{def:limite-n}
Sea $f:D\subseteq\R^n\to\R$ y $\va$ un punto de acumulación de $D$. Se dice que $\lim_{\vx\to\va}f(\vx)=L$ si para cada $\varepsilon>0$ existe $\delta>0$ tal que
\[
\vx\in D,\quad 0<\norm{\vx-\va}<\delta\ \Longrightarrow\ \abs{f(\vx)-L}<\varepsilon .
\]
\end{definicion}

La definición es la de una variable con $\abs{x-a}$ sustituido por $\norm{\vx-\va}$, y las demostraciones de las reglas algebraicas (límite de la suma, del producto, del cociente con denominador de límite no nulo) se trasladan sin cambios. Dos hechos elementales se usan constantemente: para $\vx=(x,y)$ y $\va=(a,b)$,
\begin{equation}\label{eq:t3-coordenada}
\abs{x-a}\le\norm{\vx-\va},\qquad \abs{y-b}\le\norm{\vx-\va},
\end{equation}
porque $\norm{\vx-\va}^2=(x-a)^2+(y-b)^2$ es mayor o igual que cada sumando.

\begin{proposicion}[Acotación]\label{prop:acotacion}
Si $\abs{f(\vx)-L}\le g(\vx)$ para $\vx$ cerca de $\va$ y $\lim_{\vx\to\va}g(\vx)=0$, entonces $\lim_{\vx\to\va}f(\vx)=L$.
\end{proposicion}

En el origen es útil pasar a \emph{coordenadas polares} $x=r\cos\theta$, $y=r\sin\theta$, con $r=\norm{(x,y)}$: si $\abs{f(r\cos\theta,r\sin\theta)-L}\le h(r)$ con $h(r)\to0$ cuando $r\to0^+$ y la cota $h$ \emph{no depende de $\theta$}, el límite es $L$ por la \cref{prop:acotacion}.

\addtocounter{definicion}{1}%

\begin{teorema}[Límites a lo largo de curvas]\label{teo:trayectorias}
Sea $\lim_{\vx\to\va}f(\vx)=L$ y sea $\boldsymbol\gamma:(0,\eta)\to D$ una curva con $\boldsymbol\gamma(t)\neq\va$ y $\lim_{t\to0^+}\boldsymbol\gamma(t)=\va$. Entonces $\lim_{t\to0^+}f(\boldsymbol\gamma(t))=L$.
\end{teorema}

\begin{corolario}[Criterio de no existencia]\label{cor:no-existe}
Si a lo largo de dos curvas que llegan a $\va$ la función tiene límites distintos (o a lo largo de una no tiene límite), entonces $\lim_{\vx\to\va}f(\vx)$ no existe.
\end{corolario}

\addtocounter{definicion}{1}%

\section{Continuidad}

\begin{definicion}[Continuidad]\label{def:continua-n}
Una función $f:D\to\R$ es \emph{continua en} $\va\in D$ si $\lim_{\vx\to\va}f(\vx)=f(\va)$ (o si $\va$ es un punto aislado de $D$). Es continua en $D$ si lo es en cada punto de $D$.
\end{definicion}

\begin{proposicion}[Funciones continuas elementales]\label{prop:continuas}
\begin{enumerate}
  \item Las proyecciones $(x,y)\mapsto x$ y $(x,y)\mapsto y$ son continuas en $\R^2$.
  \item Sumas, productos y cocientes (donde el denominador no se anula) de funciones continuas son continuos. En particular, los polinomios en $x$ e $y$ son continuos en $\R^2$, y las funciones racionales lo son en su dominio.
  \item Si $g$ es continua en $\va$ y $h:\R\to\R$ es continua en $g(\va)$, entonces $h\circ g$ es continua en $\va$.
\end{enumerate}
\end{proposicion}

Así, $e^{x^2+y^2}$, $\sin(xy)$ o $\ln(1+x^2+y^2)$ son continuas en todo $\R^2$, y para calcular su límite en un punto basta sustituir. Los casos interesantes son los cocientes en puntos donde se anula el denominador, como en los ejemplos anteriores.

\addtocounter{definicion}{1}%

\addtocounter{definicion}{1}%

\addtocounter{definicion}{1}%


% =====================================================================
\chapter{Derivadas parciales, gradiente y regla de la cadena}
\label{tema:4}
% =====================================================================

La derivada de una función de una variable mide cuánto cambia la función por unidad de cambio de la variable. Con varias variables hay que elegir \emph{en qué dirección} se mueve uno. Las derivadas parciales miden el cambio en las direcciones de los ejes; el gradiente las reúne en un vector que apunta hacia donde la función crece más deprisa y que es perpendicular a las curvas de nivel. Este es el objeto central del entrenamiento de redes neuronales. La presentación sigue las sesiones~26--27 y~32--38 de 18.02SC~\cite{mit1802}, las §4.3--4.6 de OpenStax~\cite{openstax3}, los capítulos~13 de Strang~\cite{strangcalc} y~8 de Apostol~\cite{apostol2}.

\section{Derivadas parciales}

\begin{definicion}[Derivadas parciales]\label{def:parcial}
Sea $f$ una función de $(x,y)$ definida cerca de $(a,b)$. Las \emph{derivadas parciales} de $f$ en $(a,b)$ son
\begin{align*}
f_x(a,b)&=\pd fx(a,b)=\lim_{h\to0}\frac{f(a+h,b)-f(a,b)}{h},\\
f_y(a,b)&=\pd fy(a,b)=\lim_{h\to0}\frac{f(a,b+h)-f(a,b)}{h},
\end{align*}
cuando existen. Para tres o más variables la definición es análoga.
\end{definicion}

Para calcular $f_x$ se deriva respecto de $x$ tratando $y$ como una constante, con todas las reglas de una variable. Geométricamente, $f_x(a,b)$ es la pendiente de la curva que se obtiene al cortar la gráfica con el plano vertical $y=b$.

\addtocounter{definicion}{1}%

La igualdad de las derivadas cruzadas del ejemplo no es casual.

\begin{teorema}[Schwarz]\label{teo:schwarz}
Si $f_{xy}$ y $f_{yx}$ existen y son continuas en un entorno de $(a,b)$, entonces $f_{xy}(a,b)=f_{yx}(a,b)$.
\end{teorema}

La demostración, que usa dos veces el teorema del valor medio, puede consultarse en Apostol~\cite[cap.~8]{apostol2}. Todas las funciones de este módulo construidas con polinomios, exponenciales, logaritmos y funciones trigonométricas cumplen la hipótesis en el interior de su dominio.

\section{Diferenciabilidad y plano tangente}

La mera existencia de las derivadas parciales es una condición débil: solo informa de lo que ocurre sobre dos rectas. La noción adecuada exige que $f$ se parezca a una función afín en \emph{todas} las direcciones.

\begin{definicion}[Diferenciabilidad]\label{def:diferenciable}
Una función $f$ definida cerca de $\va\in\R^n$ es \emph{diferenciable} en $\va$ si existe un vector $\mathbf g\in\R^n$ tal que
\[
\lim_{\mathbf h\to\vcero}\frac{f(\va+\mathbf h)-f(\va)-\mathbf g\cdot\mathbf h}{\norm{\mathbf h}}=0 .
\]
Es decir, $f(\va+\mathbf h)=f(\va)+\mathbf g\cdot\mathbf h+\varepsilon(\mathbf h)\norm{\mathbf h}$ con $\varepsilon(\mathbf h)\to0$.
\end{definicion}

\begin{teorema}\label{teo:diferenciable}
Si $f$ es diferenciable en $\va$, entonces:
\begin{enumerate}
  \item $f$ es continua en $\va$;
  \item existen las derivadas parciales en $\va$ y $\mathbf g=\nabla f(\va)$, donde el \emph{gradiente} es el vector $\nabla f(\va)=\bigl(f_{x_1}(\va),\dots,f_{x_n}(\va)\bigr)$.
\end{enumerate}
\end{teorema}

El recíproco de (ii) es falso (\cref{eje:4-no-diferenciable}). Una condición suficiente práctica es la siguiente, cuya demostración puede verse en Apostol~\cite[cap.~8]{apostol2}: \emph{si las derivadas parciales de $f$ existen y son continuas en un entorno de $\va$, entonces $f$ es diferenciable en $\va$.}

\begin{definicion}[Plano tangente y aproximación lineal]\label{def:tangente}
Si $f$ es diferenciable en $(a,b)$, el \emph{plano tangente} a la gráfica de $f$ en $(a,b,f(a,b))$ es
\[
z=f(a,b)+f_x(a,b)(x-a)+f_y(a,b)(y-b),
\]
y la \emph{aproximación lineal} es $f(a+\Delta x,b+\Delta y)\approx f(a,b)+f_x(a,b)\Delta x+f_y(a,b)\Delta y$. Con notación de diferenciales, $df=f_x\,dx+f_y\,dy$.
\end{definicion}

\addtocounter{definicion}{1}%

\section{Regla de la cadena}

\begin{teorema}[Regla de la cadena]\label{teo:cadena}
Sea $\vr(t)=(x(t),y(t))$ derivable en $t_0$ y $f$ diferenciable en $\vr(t_0)$. Entonces $w(t)=f(\vr(t))$ es derivable en $t_0$ y
\[
\frac{dw}{dt}=f_x\frac{dx}{dt}+f_y\frac{dy}{dt}=\nabla f(\vr(t_0))\cdot\vr\,'(t_0).
\]
\end{teorema}

Si $x$ e $y$ dependen de dos variables $(u,v)$, se aplica el teorema fijando una de ellas:
\begin{equation}\label{eq:t4-cadena2}
\pd wu=\pd wx\pd xu+\pd wy\pd yu,\qquad \pd wv=\pd wx\pd xv+\pd wy\pd yv .
\end{equation}
La regla se recuerda con un \emph{diagrama de árbol}: cada camino de $w$ a la variable final aporta el producto de las derivadas de sus tramos, y se suman las aportaciones de todos los caminos.

\addtocounter{figure}{1}%

\addtocounter{definicion}{1}%

\section{Gradiente y derivada direccional}

\begin{definicion}[Derivada direccional]\label{def:direccional}
Sea $\vu$ un vector \emph{unitario}. La \emph{derivada direccional} de $f$ en $\va$ según $\vu$ es
\[
D_{\vu}f(\va)=\lim_{t\to0}\frac{f(\va+t\vu)-f(\va)}{t},
\]
la tasa de cambio de $f$ por unidad de longitud recorrida en la dirección de $\vu$. Para $\vu=\vi$ y $\vu=\vj$ se obtienen $f_x$ y $f_y$.
\end{definicion}

\begin{teorema}[Cálculo de la derivada direccional]\label{teo:direccional}
Si $f$ es diferenciable en $\va$, entonces $D_{\vu}f(\va)=\nabla f(\va)\cdot\vu$ para todo vector unitario $\vu$.
\end{teorema}

\begin{corolario}[Dirección de máximo crecimiento]\label{cor:maximo-crecimiento}
Sea $\nabla f(\va)\neq\vcero$. Entre todas las direcciones unitarias, $D_{\vu}f(\va)$ alcanza su máximo, $\norm{\nabla f(\va)}$, para $\vu=\nabla f(\va)/\norm{\nabla f(\va)}$, su mínimo, $-\norm{\nabla f(\va)}$, en la dirección opuesta, y se anula en las direcciones perpendiculares al gradiente.
\end{corolario}

\begin{teorema}[El gradiente es normal a los conjuntos de nivel]\label{teo:gradiente-normal}
Sea $f$ diferenciable y $\vr(t)$ una curva derivable contenida en el conjunto de nivel $f=c$, con $\vr(t_0)=\va$. Entonces $\nabla f(\va)\cdot\vr\,'(t_0)=0$.
\end{teorema}

En consecuencia, el \emph{plano tangente} a la superficie de nivel $F(x,y,z)=c$ en un punto $P_0$ con $\nabla F(P_0)\neq\vcero$ es el plano de normal $\nabla F(P_0)$: $\nabla F(P_0)\cdot(X-P_0)=0$. El plano tangente a una gráfica $z=f(x,y)$ es el caso particular $F(x,y,z)=f(x,y)-z$.

\addtocounter{definicion}{1}%

\addtocounter{definicion}{1}%

\addtocounter{definicion}{1}%


% =====================================================================
\chapter{Extremos libres y condicionados}
\label{tema:5}
% =====================================================================

Entrenar un modelo es minimizar una función de muchas variables, y estimar probabilidades por máxima verosimilitud es maximizar una función con la restricción de que las probabilidades sumen uno. Este tema estudia ambos problemas en dimensión dos y tres: los extremos libres, que se localizan donde se anula el gradiente y se clasifican con la matriz hessiana, y los extremos condicionados, que se resuelven con multiplicadores de Lagrange. La presentación sigue las sesiones~28--31 y~39--40 de 18.02SC~\cite{mit1802}, las §4.7--4.8 de OpenStax~\cite{openstax3}, las §9.9--9.12 y~9.14 de Apostol~\cite{apostol2}, el texto de Marsden y Tromba~\cite{marsden} y las §13.6--13.7 de Strang~\cite{strangcalc}.

\section{Extremos locales y puntos críticos}

\begin{definicion}[Extremos]\label{def:extremos}
Sea $f:D\subseteq\R^n\to\R$ y $\va\in D$. Se dice que $f$ tiene en $\va$ un \emph{máximo local} si existe $\delta>0$ tal que $f(\vx)\le f(\va)$ para todo $\vx\in D$ con $\norm{\vx-\va}<\delta$; un \emph{mínimo local}, con $\ge$. Si la desigualdad vale para todo $\vx\in D$, el extremo es \emph{absoluto} (o global). Un punto $\va$ interior al dominio es \emph{crítico} si $\nabla f(\va)=\vcero$ o si alguna derivada parcial no existe en $\va$.
\end{definicion}

\begin{teorema}[Condición necesaria de Fermat]\label{teo:fermat}
Si $f$ tiene un extremo local en un punto interior $\va$ y sus derivadas parciales existen en $\va$, entonces $\nabla f(\va)=\vcero$.
\end{teorema}

El recíproco es falso: $f(x,y)=x^2-y^2$ tiene $\nabla f(0,0)=\vcero$ y no tiene extremo en el origen, pues crece a lo largo del eje $x$ y decrece a lo largo del eje $y$. Un punto crítico que no es extremo local se llama \emph{punto de silla}.

\section{Clasificación: el criterio del hessiano}

\begin{definicion}[Matriz hessiana]\label{def:hessiana}
Si $f(x,y)$ tiene derivadas segundas, su \emph{matriz hessiana} y su \emph{discriminante} en $\va$ son
\[
H_f(\va)=\begin{pmatrix}f_{xx}&f_{xy}\\f_{yx}&f_{yy}\end{pmatrix}(\va),\qquad
\Delta(\va)=\det H_f(\va)=f_{xx}f_{yy}-f_{xy}^2 ,
\]
donde se ha usado el teorema de Schwarz para escribir $f_{xy}=f_{yx}$.
\end{definicion}

La idea del criterio es la fórmula de Taylor de segundo orden: si $f$ tiene derivadas segundas continuas y $\nabla f(\va)=\vcero$, entonces, con $\mathbf h=(h,k)$,
\begin{equation}\label{eq:t5-taylor}
f(\va+\mathbf h)=f(\va)+\tfrac12\,Q(\mathbf h)+\varepsilon(\mathbf h)\norm{\mathbf h}^2,\qquad Q(h,k)=Ah^2+2Bhk+Ck^2,
\end{equation}
con $A=f_{xx}(\va)$, $B=f_{xy}(\va)$, $C=f_{yy}(\va)$ y $\varepsilon(\mathbf h)\to0$. La demostración de~\eqref{eq:t5-taylor} aplica la fórmula de Taylor de una variable a $t\mapsto f(\va+t\mathbf h)$ (Apostol~\cite[§9.9--9.12]{apostol2}; módulo~4, tema~4). Cerca de un punto crítico, por tanto, el signo de $f(\va+\mathbf h)-f(\va)$ lo decide la \emph{forma cuadrática} $Q$.

\begin{lema}[Signo de una forma cuadrática]\label{lem:forma}
Sea $Q(h,k)=Ah^2+2Bhk+Ck^2$ y $\Delta=AC-B^2$.
\begin{enumerate}
  \item Si $\Delta>0$ y $A>0$, entonces $Q(h,k)>0$ para todo $(h,k)\neq(0,0)$.
  \item Si $\Delta>0$ y $A<0$, entonces $Q(h,k)<0$ para todo $(h,k)\neq(0,0)$.
  \item Si $\Delta<0$, $Q$ toma valores positivos y negativos.
\end{enumerate}
\end{lema}

\begin{teorema}[Criterio del hessiano]\label{teo:hessiano}
Sea $f$ con derivadas segundas continuas cerca de un punto crítico $\va$ con $\nabla f(\va)=\vcero$, y sea $\Delta=\Delta(\va)$.
\begin{enumerate}
  \item Si $\Delta>0$ y $f_{xx}(\va)>0$, $f$ tiene un mínimo local en $\va$.
  \item Si $\Delta>0$ y $f_{xx}(\va)<0$, $f$ tiene un máximo local en $\va$.
  \item Si $\Delta<0$, $\va$ es un punto de silla.
  \item Si $\Delta=0$, el criterio no decide.
\end{enumerate}
\end{teorema}

\addtocounter{definicion}{1}%

\section{Extremos absolutos en regiones cerradas y acotadas}

\begin{teorema}[Weierstrass]\label{teo:weierstrass}
Una función continua sobre un conjunto cerrado y acotado de $\R^n$ alcanza en él un máximo y un mínimo absolutos.
\end{teorema}

La demostración usa la completitud de los números reales y se omite; para una variable es el teorema de los valores extremos del cálculo elemental. En la práctica, para hallar los extremos absolutos de $f$ en una región $R$ cerrada y acotada: (1) se hallan los puntos críticos interiores; (2) se estudia $f$ sobre la frontera, que suele reducirse a funciones de una variable o a un problema con multiplicadores de Lagrange; (3) se comparan todos los valores obtenidos.

\addtocounter{definicion}{1}%

\section{Multiplicadores de Lagrange}

Se busca el máximo o el mínimo de $f(x,y)$ cuando $(x,y)$ está obligado a moverse sobre la curva $g(x,y)=c$. En un extremo, la curva de nivel de $f$ que pasa por el punto no puede cruzar la restricción, porque entonces se podría seguir avanzando sobre $g=c$ hacia valores mayores y menores de $f$: ambas curvas deben ser \emph{tangentes}, y sus gradientes, paralelos (\cref{fig:t5-lagrange}).

\addtocounter{figure}{1}%

\begin{teorema}[Multiplicadores de Lagrange]\label{teo:lagrange}
Sean $f$ y $g$ funciones de $(x,y)$ con derivadas parciales continuas, y sea $\va$ un punto de la curva $g=c$ con $\nabla g(\va)\neq\vcero$. Si $f$, restringida a la curva, tiene un extremo local en $\va$, existe $\lambda\in\R$ (el \emph{multiplicador}) tal que
\[
\nabla f(\va)=\lambda\nabla g(\va).
\]
El mismo enunciado vale para $f$ y $g$ de tres variables y la superficie $g=c$.
\end{teorema}

El método consiste en resolver el sistema formado por $\nabla f=\lambda\nabla g$ y $g=c$ (tres ecuaciones con incógnitas $x,y,\lambda$ en el plano; cuatro en el espacio). Sus soluciones son los \emph{candidatos}; si la restricción es cerrada y acotada, el teorema de Weierstrass garantiza que el mayor valor de $f$ entre ellos es el máximo y el menor, el mínimo.

\addtocounter{definicion}{1}%

\addtocounter{definicion}{1}%

\addtocounter{definicion}{1}%

\addtocounter{definicion}{1}%


% =====================================================================
\chapter{Integrales dobles y la integral gaussiana}
\label{tema:6}
% =====================================================================

La integral de una variable suma los valores de una función a lo largo de un intervalo; la integral doble los suma sobre una región del plano y mide, por ejemplo, el volumen bajo una superficie. Su cálculo se reduce a dos integrales de una variable sucesivas (teorema de Fubini), y en regiones circulares conviene usar coordenadas polares. La aplicación más importante para el curso es el cálculo de la integral gaussiana $\int_{-\infty}^{\infty}e^{-x^2}dx=\sqrt\pi$, que no puede obtenerse con primitivas elementales y sí con un rodeo por el plano. La presentación sigue las sesiones~47--54 de 18.02SC~\cite{mit1802}, las §5.1--5.3 de OpenStax~\cite{openstax3}, las §11.1--11.14 de Apostol~\cite{apostol2} y las §14.1--14.2 de Strang~\cite{strangcalc}.

\section{Integral doble sobre un rectángulo}

Sea $R=[a,b]\times[c,d]$ y $f$ acotada en $R$. Se divide $[a,b]$ en $m$ subintervalos de longitud $\Delta x$ y $[c,d]$ en $n$ de longitud $\Delta y$, lo que divide $R$ en $mn$ rectángulos $R_{ij}$ de área $\Delta A=\Delta x\,\Delta y$, y se elige un punto $(x_{ij}^*,y_{ij}^*)$ en cada uno.

\begin{definicion}[Integral doble]\label{def:integral-doble}
La \emph{integral doble} de $f$ sobre $R$ es
\[
\iint_R f(x,y)\dA=\lim_{m,n\to\infty}\sum_{i=1}^{m}\sum_{j=1}^{n}f(x_{ij}^*,y_{ij}^*)\,\Delta A ,
\]
si el límite existe y no depende de la elección de los puntos; se dice entonces que $f$ es \emph{integrable} en $R$. Toda función continua en $R$ es integrable.
\end{definicion}

Si $f\ge0$, cada sumando es el volumen de un prisma de base $R_{ij}$ y altura $f(x_{ij}^*,y_{ij}^*)$, y la integral es el \emph{volumen} del sólido comprendido entre $R$ y la gráfica de $f$. Si $f=1$, la integral es el área de $R$.

\begin{proposicion}[Propiedades]\label{prop:propiedades-integral}
Para $f,g$ integrables en $R$ y $\alpha,\beta\in\R$:
\begin{enumerate}
  \item linealidad: $\iint_R(\alpha f+\beta g)\dA=\alpha\iint_Rf\dA+\beta\iint_Rg\dA$;
  \item monotonía: si $f\le g$ en $R$, entonces $\iint_Rf\dA\le\iint_Rg\dA$;
  \item aditividad: si $R$ se divide en dos regiones $R_1$, $R_2$ que solo comparten frontera, $\iint_R f=\iint_{R_1}f+\iint_{R_2}f$.
\end{enumerate}
\end{proposicion}

Las tres propiedades se heredan de las sumas de Riemann al pasar al límite: por ejemplo, si $f\le g$, cada suma de $f$ es menor o igual que la correspondiente de $g$. El cálculo efectivo se apoya en el teorema siguiente, cuya demostración puede verse en Apostol~\cite[§11.10--11.11]{apostol2}.

\begin{teorema}[Fubini, para rectángulos]\label{teo:fubini}
Si $f$ es continua en $R=[a,b]\times[c,d]$, entonces
\[
\iint_R f\dA=\int_a^b\left(\int_c^d f(x,y)\,dy\right)dx=\int_c^d\left(\int_a^b f(x,y)\,dx\right)dy .
\]
\end{teorema}

En la integral interior $\int_c^df(x,y)\,dy$, la variable $x$ se trata como constante. Geométricamente, es el área de la sección del sólido por el plano vertical de abscisa $x$, y el volumen se obtiene sumando secciones.

\begin{corolario}[Integrandos separables]\label{cor:separable}
Si $f(x,y)=g(x)h(y)$ con $g$ y $h$ continuas, entonces $\iint_{[a,b]\times[c,d]}g(x)h(y)\dA=\left(\int_a^bg\right)\left(\int_c^dh\right)$.
\end{corolario}

\addtocounter{definicion}{1}%

\section{Regiones generales y cambio del orden de integración}

\begin{definicion}[Regiones elementales]\label{def:regiones}
Una región es \emph{de tipo~I} si es de la forma $\{a\le x\le b,\ g_1(x)\le y\le g_2(x)\}$, y \emph{de tipo~II} si es $\{c\le y\le d,\ h_1(y)\le x\le h_2(y)\}$, con $g_i$, $h_i$ continuas.
\end{definicion}

Para $f$ continua, el teorema de Fubini se extiende a estas regiones~\cite[§11.12--11.14]{apostol2}:
\begin{align*}
\iint_Df\dA&=\int_a^b\int_{g_1(x)}^{g_2(x)}f(x,y)\,dy\,dx&&\text{(tipo I)},\\
\iint_Df\dA&=\int_c^d\int_{h_1(y)}^{h_2(y)}f(x,y)\,dx\,dy&&\text{(tipo II)}.
\end{align*}
Los límites de la integral interior pueden depender de la variable exterior; los de la exterior han de ser constantes. El método práctico es: dibujar la región, fijar la variable exterior y leer, sobre un segmento vertical (tipo~I) u horizontal (tipo~II), dónde entra y dónde sale de la región.

\addtocounter{definicion}{1}%

Cuando la integral interior no tiene primitiva elemental, puede bastar con cambiar el orden. Hay que describir la \emph{misma} región del otro modo.

\addtocounter{definicion}{1}%

\section{Coordenadas polares}

En polares, $x=r\cos\theta$, $y=r\sin\theta$, con $r\ge0$. Las regiones naturales son los \emph{rectángulos polares} $\{r_1\le r\le r_2,\ \theta_1\le\theta\le\theta_2\}$, sectores de corona circular. Su área es la diferencia de dos sectores circulares:
\[
\tfrac12r_2^2\Delta\theta-\tfrac12r_1^2\Delta\theta=\tfrac{r_1+r_2}{2}\,(r_2-r_1)\,\Delta\theta=\bar r\,\Delta r\,\Delta\theta ,
\]
donde $\bar r$ es el radio medio. Así, un rectángulo polar pequeño tiene área aproximada $r\,\Delta r\,\Delta\theta$, no $\Delta r\,\Delta\theta$: los rectángulos más alejados del origen son más grandes. Sumando y pasando al límite se obtiene el resultado siguiente~\cite[§11.26--11.31]{apostol2}.

\begin{teorema}[Cambio a coordenadas polares]\label{teo:polares}
Si $f$ es continua en la región $D=\{\alpha\le\theta\le\beta,\ r_1(\theta)\le r\le r_2(\theta)\}$, con $0\le\beta-\alpha\le2\pi$, entonces
\[
\iint_Df(x,y)\dA=\int_\alpha^\beta\int_{r_1(\theta)}^{r_2(\theta)}f(r\cos\theta,r\sin\theta)\,r\,dr\,d\theta .
\]
\end{teorema}

El factor $r$ es el valor absoluto del determinante de las derivadas parciales de $(x,y)$ respecto de $(r,\theta)$, el \emph{jacobiano} $\bigl|\begin{smallmatrix}\cos\theta&-r\sin\theta\\\sin\theta&r\cos\theta\end{smallmatrix}\bigr|=r$, de acuerdo con la interpretación del determinante como factor de cambio de área (\cref{teo:det-area}).

\addtocounter{definicion}{1}%

\section{La integral gaussiana}

\begin{teorema}[Integral gaussiana]\label{teo:gauss}
La integral impropia $\int_{-\infty}^{\infty}e^{-x^2}dx$ converge y vale $\sqrt\pi$.
\end{teorema}

\addtocounter{figure}{1}%

\begin{corolario}\label{cor:gauss-normal}
$\displaystyle\int_{-\infty}^{\infty}e^{-x^2/2}\,dx=\sqrt{2\pi}$. En consecuencia, la densidad normal estándar $\varphi(x)=\frac{1}{\sqrt{2\pi}}e^{-x^2/2}$ integra~$1$.
\end{corolario}

\addtocounter{definicion}{1}%

\addtocounter{definicion}{1}%



\backmatter
% =====================================================================
\begin{thebibliography}{99}
\addcontentsline{toc}{chapter}{Bibliografía}
\markboth{Bibliografía}{}
\small

\bibitem{apostol2}
T.~M. Apostol.
\newblock \emph{Calculus}, vol.~2: \emph{Multi-Variable Calculus and Linear Algebra, with Applications to Differential Equations and Probability}, 2.ª ed.
\newblock Wiley, 1969.
\newblock Caps.~8 «Differential Calculus of Scalar and Vector Fields», 9 «Applications of Differential Calculus» y 11 «Multiple Integrals». Texto del curso 18.024 del MIT.

\bibitem{mit1802}
D.~Auroux \emph{et al.}
\newblock 18.02SC: Multivariable Calculus.
\newblock Massachusetts Institute of Technology, MIT OpenCourseWare, otoño de 2010.
\newblock Unidades~1 «Vectors and Matrices», 2 «Partial Derivatives» y 3 «Double Integrals and Line Integrals in the Plane»; \emph{problem sets} 1--7 y problemas suplementarios (secciones 1A--1H, 2A--2J y 3A--3D).
\newblock \url{https://ocw.mit.edu/courses/18-02sc-multivariable-calculus-fall-2010/}

\bibitem{openstax3}
G.~Strang y E.~J. Herman (autores principales).
\newblock \emph{Calculus}, vol.~3.
\newblock OpenStax, Rice University, 2016.
\newblock Caps.~2 «Vectors in Space», 4 «Differentiation of Functions of Several Variables» y 5 «Multiple Integration».
\newblock \url{https://openstax.org/details/books/calculus-volume-3}

\bibitem{marsden}
J.~E. Marsden y A.~Tromba.
\newblock \emph{Vector Calculus}, 6.ª ed.
\newblock W.~H. Freeman, 2012.

\bibitem{strangla}
G.~Strang.
\newblock \emph{Introduction to Linear Algebra}, 6.ª ed.
\newblock Wellesley-Cambridge Press, 2023.
\newblock §2.2 «Elimination Matrices and Inverse Matrices». Texto del curso 18.06 del MIT.
\newblock \url{https://math.mit.edu/~gs/linearalgebra/ila6/indexila6.html}

\bibitem{strangcalc}
G.~Strang.
\newblock \emph{Calculus}, 3.ª ed. actualizada.
\newblock Wellesley-Cambridge Press. Edición libre en MIT OpenCourseWare (RES.18-001, otoño de 2023).
\newblock Caps.~11 «Vectors and Matrices», 13 «Partial Derivatives» y 14 «Multiple Integrals».
\newblock \url{https://ocw.mit.edu/courses/res-18-001-calculus-fall-2023/}

\end{thebibliography}


\end{document}
